\documentclass[11pt]{article}
\usepackage[margin=1in]{geometry}
\usepackage{amsmath,amssymb,amsthm}
\usepackage{mathtools}
\usepackage{booktabs}
\usepackage{array}
\usepackage{enumitem}
\usepackage{tabularx}
\usepackage[numbers,sort]{natbib}
\PassOptionsToPackage{obeyspaces,spaces}{url}
\usepackage[hidelinks]{hyperref}
\usepackage{url}

% Bold math in bold headings (article defaults plus \boldmath).
\makeatletter
\renewcommand\section{\@startsection{section}{1}{\z@}%
  {-3.5ex \@plus -1ex \@minus -.2ex}{2.3ex \@plus.2ex}%
  {\normalfont\Large\bfseries\boldmath}}
\renewcommand\subsection{\@startsection{subsection}{2}{\z@}%
  {-3.25ex\@plus -1ex \@minus -.2ex}{1.5ex \@plus .2ex}%
  {\normalfont\large\bfseries\boldmath}}
\renewcommand\paragraph{\@startsection{paragraph}{4}{\z@}%
  {3.25ex \@plus1ex \@minus.2ex}{-1em}%
  {\normalfont\normalsize\bfseries\boldmath}}
\makeatother

% Lists: a small gap above and below, none between items.
\setlist[itemize]{topsep=0.4\baselineskip,itemsep=0pt,parsep=0pt}

% Notation
\newcommand{\B}{\{0,1\}}
\newcommand{\Bs}{\{0,1\}^*}
\newcommand{\eps}{\epsilon}
\newcommand{\E}{\mathbb{E}}
\newcommand{\N}{\mathbb{N}}
\newcommand{\cM}{\mathcal{M}}
\newcommand{\cMone}{\mathcal{M}_1}
\newcommand{\cCQ}{\mathcal{C}_{\mathbb{Q}}}
\newcommand{\cCcomp}{\mathcal{C}_{\mathrm{comp}}}
\newcommand{\ind}{\mathbf{1}}
\newcommand{\Mnorm}{M_{\mathrm{norm}}}
\newcommand{\geqx}{\mathrel{\stackrel{\times}{\geq}}}
\newcommand{\leqx}{\mathrel{\stackrel{\times}{\leq}}}
\newcommand{\eqx}{\mathrel{\stackrel{\times}{=}}}
\newcommand{\KL}{D_{\mathrm{KL}}}
\newcommand{\AI}{\mathrm{AI}}
\newcommand{\SP}{\mathrm{SP}}
\newcommand{\KM}{\mathrm{KM}}
\newcommand{\AIxi}{\textup{AI}\ensuremath{\xi}}
\newcommand{\lean}[1]{\path{#1}}
\newcommand{\todo}[1]{\textbf{[TODO: #1]}}
% Replace with the public repository URL, e.g. \url{https://...}.
\newcommand{\repourl}{\url{https://github.com/mattmotoki/uai-book-problems}}

\newtheorem{theorem}{Theorem}[section]
\newtheorem{lemma}[theorem]{Lemma}
\newtheorem{proposition}[theorem]{Proposition}
\newtheorem{corollary}[theorem]{Corollary}
\theoremstyle{definition}
\newtheorem{definition}[theorem]{Definition}
\theoremstyle{remark}
\newtheorem{remark}[theorem]{Remark}

\title{Answers to three prize problems from Hutter's\\ \emph{Universal Artificial Intelligence}}
\author{Matt Motoki\\
Independent researcher\\
\texttt{mattmotoki@gmail.com}}
\date{Draft, November 2026}

\begin{document}
\maketitle

\begin{abstract}
We report answers to three prize problems from Marcus Hutter's book
\emph{Universal Artificial Intelligence} (2005).
\emph{Problem 3.3} asks whether the cumulative 0-1 losses of the Bayes
predictors of two mixtures over the same class differ by a finite amount.
The answer is no, also in the book's own setting, where the class is a
countable class of probability measures. For the class of probability
measures with computable rational conditionals, and for every positive prior
$a$ on it, we give two priors $w^\pm\ge2^{-20}a$ whose ratio is between $1/3$
and $3$, and a computable measure $\mu$ in the class that does not depend on
$a$, such that the $\mu$-expected loss difference of the two Bayes mixtures is
at least $n^{1/4}/256$ after $n\ge2$ steps. A paper step extends this to all
computable probability measures. The same bound holds for two Bayes mixtures
over an enumeration of all lower semicomputable semimeasures, also with the
exact weights $2^{-K(i)}$ of the universal prior, for two constructed optimal
prefix machines and one constructed enumeration, where the two mixtures differ
only in the machine that defines $K$ (formal bound $n^{1/4}/512$ for $n\ge3$). A simpler two-measure example has a difference of
at least $H_n/32$. For log loss, quadratic loss, and sources with a uniform
decision margin, the difference is known to stay bounded.
\emph{Problem 6.3} ends with the question whether the number of errors of the
universal agent in cycles $k$ to $n$ is bounded by a constant times
$2^{K(\dot z_{k:n}\mid \dot z_{<k})}$. For ordinary conditional prefix
complexity, the answer is no, for every deterministic predictor.
\emph{Problem 2.7} asks whether the normalized universal predictor $\Mnorm$
has lower error bounds of order $2^{-K(t)}$. We prove that, for every
universal semimeasure and every computable measure, the expected one-step error
of $\Mnorm$ is at least $c\,2^{-K(t)}$ at some time in every late window of
length $O(\log t)$, and along every computable selection of times. We also
show that no positive power of $2^{-K(t)}$ is an upper bound for all universal
$\Mnorm$, already for the fair coin. The bound at every large $t$ stays open;
we show that a set of times along which the expected error divided by
$2^{-K(t)}$ tends to $0$ must be immune.
The main theorems are checked in Lean~4.
\end{abstract}

\tableofcontents

%======================================================================
\section{Introduction}
%======================================================================

Hutter's monograph \emph{Universal Artificial Intelligence}~\cite{hutter2005uai}
ends each chapter with problems, and its web page~\cite{hutterprizepage} offers
prizes for some of them. This paper reports our answers to three of them.

\paragraph{Problem 3.3 (comparing two mixtures).}
Let $\xi$ and $\xi'$ be two Bayes mixtures over the same class $\cM$, with
positive prior weights, and let the true environment $\mu$ be in $\cM$. The
known part of the problem bounds the $\mu$-expected squared distance between
their predictions and the difference of their losses by
$O\bigl((L^{\Lambda_\mu}_n)^{1/2}\bigr)$. The book then asks~\cite[p.~120]{hutter2005uai}:
\begin{quote}
``Is the stronger result $L_\infty^{\Lambda_\xi}-L_\infty^{\Lambda_{\xi'}}<\infty$ also true?''
\end{quote}
Both total losses are usually infinite, so we read the question as:\ is the
sequence $L_n^{\Lambda_\xi}-L_n^{\Lambda_{\xi'}}$ bounded? We answer no
(Section~\ref{sec:33}). In the book's Section~3.2.2, where the problem is set,
$\cM$ is a countable class of probability measures~\cite[p.~70]{hutter2005uai}.
Theorem~\ref{thm:measures} answers the problem in this setting. Take the class
of probability measures with computable rational conditionals, and any prior
$a$ on it with $a_\nu>0$ and $\sum_\nu a_\nu=1$. Then there are priors
$w^\pm\ge2^{-20}a$ with $w^-_\nu/w^+_\nu\in[1/3,3]$ whose Bayes mixtures have a
loss difference of at least $n^{1/4}/256$ under one fixed computable $\mu$ in
the class. A paper step (Corollary~\ref{cor:allmeasures}) gives the same result
for all computable probability measures. The same construction gives two Bayes
mixtures over an enumeration of all lower semicomputable semimeasures
(Theorem~\ref{thm:mixture}), with weights $w^\pm_i>0$, $\sum_i w^\pm_i=1$ and
$w^-_i/w^+_i\in[1/3,3]$, a version with weights within constant factors of
$2^{-K(i)}$ (Theorem~\ref{thm:complexity}), and a version with the exact
weights: two universal mixtures $\sum_i2^{-K_C(i)}\tau_i$ over one constructed
enumeration $\tau$ that differ only in the constructed optimal prefix machine $C$
(Theorem~\ref{thm:exact}). All three have the bound $n^{1/4}/256$ for $n\ge2$
(for Theorem~\ref{thm:exact}, the formal bound is $n^{1/4}/512$ for $n\ge3$).
Section~\ref{sec:33} also marks where the answer turns positive: log loss,
quadratic loss, and sources with a uniform decision margin.

\paragraph{Problem 6.3 (posterization of prediction errors).}
The problem leads to a bound on the number $E_{k:n}$ of errors of the agent
\AIxi{} in cycles $k$ to $n$,
\[
E_{k:n}\leqx\xi^{\AI}(\dot y\dot x_{<k})/\xi^{\SP}(\dot z_{1:n}),
\]
where $\dot y\dot x_{<k}$ is the action-perception history before cycle $k$,
and then asks~\cite[p.~208]{hutter2005uai}:
\begin{quote}
``Is it possible to improve the numerator to $\xi^{\SP}(\dot z_{<k})$ and to
bound the expression by $\approx 2^{K(\dot z_{k:n}|\dot z_{<k})}$ (cf.\ Problem 3.13)?''
\end{quote}
The two improvements form a chain
$E_{k:n}\lesssim \xi^{\SP}(\dot z_{<k})/\xi^{\SP}(\dot z_{1:n})\lesssim 2^{K(\dot z_{k:n}|\dot z_{<k})}$.
We show that the end of the chain is false: with ordinary conditional prefix
complexity, no constant $C$ gives
$E_{k:n}\le C\,2^{K(\dot z_{k:n}|\dot z_{<k})}$, for \AIxi{} with any fixed tie
rule and for every other deterministic predictor
(Theorem~\ref{thm:63E}, Section~\ref{sec:63}). The first link of the chain and
the version with Chernov and Hutter's complexity $K_*$ stay open.

\paragraph{Problem 2.7 (lower convergence bounds and defect of $M$).}
For Solomonoff's universal semimeasure $M$, the book proves a lower bound of
order $2^{-K(t)}$ on the one-step prediction error at every time $t$; the bound
comes from the defect $1-\sum_{x_t}M(x_t|x_{<t})$. It then
asks~\cite[p.~62]{hutter2005uai}:
\begin{quote}
``Do similar lower performance bounds hold for measure
\[\Mnorm(x_t|x_{<t}):=M(x_t|x_{<t})/\textstyle\sum_{x_t}M(x_t|x_{<t}),\]
which has no defect? For which sequences does equality $\eqx 2^{-K(t)}$ hold?''
\end{quote}
For deterministic environments the answer follows from Hutter's later
paper~\cite{hutter2007bayesconf} (Proposition~\ref{prop:27F}). For stochastic
environments we prove expected lower bounds of order $2^{-K(t)}$ at a time in
every late window of length $O(\log t)$, and infinitely often along every
computable selection of times, with a constant that depends on the selection
(Theorem~\ref{thm:selection} and its corollaries). On the upper side, we show that some universal $M$ has fair-coin
expected errors that are not $O(2^{-\alpha K(t)})$ for any $\alpha>0$
(Theorem~\ref{thm:spikes}). Whether the lower bound holds at every large $t$
stays open; Proposition~\ref{prop:immune} gives a constraint on a
counterexample.

\paragraph{Method and formal verification.}
For each result, we first fixed a Lean~4 statement and had it reviewed against
the book's text. The proofs were then checked in Lean~4 with
Mathlib~\cite{mathlib2020} and the library
\texttt{kolmogorov-complexity-lean}~\cite{milovanovlean}.
Section~\ref{sec:formal} lists the theorem names. The paper proofs below are
complete and do not depend on the formal proofs. Where the formal statement is
narrower than the paper statement (rational conditional probabilities, a class
indexed by its conditionals, semimeasures with root mass $1$, a fixed tie
rule, a weaker constant), we say so at the theorem.

\paragraph{Summary.}
Table~\ref{tab:summary} lists the results.

\begin{table}[ht]
\centering\small
\begin{tabularx}{\textwidth}{@{}>{\raggedright\arraybackslash}p{3.9cm}>{\raggedright\arraybackslash}X>{\raggedright\arraybackslash}p{3.7cm}@{}}
\toprule
Result & Statement (short) & Status \\
\midrule
Thm.~\ref{thm:warmup} & 3.3, two measures: difference $\ge H_n/32$ & new answer; Lean \\
\addlinespace
Thm.~\ref{thm:robust}, \ref{thm:aggregate} & finite robust experiment; aggregation over any positive background & Lean \\
\addlinespace
Thm.~\ref{thm:mixture} & 3.3, two Bayes mixtures over all lsc semimeasures: difference $\ge n^{1/4}/256$ & new answer; Lean \\
\addlinespace
Thm.~\ref{thm:complexity} & 3.3 with weights $\asymp 2^{-K(i)}$ & new answer; Lean \\
\addlinespace
Thm.~\ref{thm:exact} & 3.3 with exact weights $2^{-K_{C_\pm}(i)}$: two universal mixtures that differ only in the constructed optimal prefix machines $C_\pm$ & new answer; Lean \\
\addlinespace
Thm.~\ref{thm:measures}, Cor.~\ref{cor:allmeasures} & 3.3 in the book's setting: two Bayes mixtures over the probability measures with computable rational conditionals (Cor.: all computable probability measures), priors $w^\pm\ge2^{-20}a$ from any positive base prior $a$: difference $\ge n^{1/4}/256$ & new answer; Lean (Cor.: paper step) \\
\addlinespace
Thm.~\ref{thm:63E} & 6.3, last question: no bound $C\,2^{K(y|x)}$ & new answer (extends a known mechanism); Lean \\
\addlinespace
Prop.~\ref{prop:27F} & 2.7, deterministic: $\Mnorm(b|x)\ge c\,2^{-K(t)}$ & known~\cite{hutter2007bayesconf}; Lean for one $M$ \\
\addlinespace
Thm.~\ref{thm:selection}, Cor.~\ref{cor:window}, \ref{cor:dyadic} & 2.7, expected lower bounds along selections, in $O(\log t)$ windows, in dyadic blocks & partial answer; Lean \\
\addlinespace
Thm.~\ref{thm:spikes} & 2.7, no upper bound $O(2^{-\alpha K(t)})$ for all $M$ & partial answer; Lean \\
\addlinespace
Prop.~\ref{prop:immune} & 2.7, quiet sets of times (where error$/2^{-K(t)}\to0$) are immune & constraint on the open case \\
\bottomrule
\end{tabularx}
\caption{Results. ``Lean'' means that the stated formal version has a complete
Lean proof that uses only the three standard axioms.}
\label{tab:summary}
\end{table}

%======================================================================
\section{Preliminaries}\label{sec:prelim}
%======================================================================

\paragraph{Strings.}
$\B^*$ is the set of finite binary strings, $\eps$ the empty string, $|x|$ the
length of $x$. For a string or sequence $x$, $x_{<t}=x_1\dots x_{t-1}$ and
$x_{k:n}=x_k\dots x_n$. We write $\bar b=1-b$ for a bit $b$.

\paragraph{Semimeasures.}
A \emph{semimeasure} is a function $\nu:\Bs\to[0,1]$ with $\nu(\eps)\le1$ and
$\nu(x0)+\nu(x1)\le\nu(x)$ for all $x$. It is a \emph{(probability) measure}
if $\nu(\eps)=1$ and equality holds for all $x$. A function
$f:\Bs\to[0,\infty)$ is \emph{lower semicomputable} (lsc) if there is a
computable $\phi:\Bs\times\N\to\mathbb{Q}$, nondecreasing in its second
argument, with $\lim_s\phi(x,s)=f(x)$. A family is uniformly lsc if one
$\phi$ works for all members. Let $\cM$ be the set of lsc semimeasures (the
book's class $\cM_U$, \cite[p.~48]{hutter2005uai}) and $\cMone\subseteq\cM$
the members with $\nu(\eps)=1$. A measure $\mu$ is \emph{computable} if
$\mu(x)$ is a computable real, uniformly in $x$; it has \emph{rational
conditionals} if $x\mapsto\mu(1|x)$ is a computable function into
$\mathbb Q$. Computable measures are lsc, so they belong to $\cMone$.
A semimeasure $M$ is \emph{universal} if it
is lsc and for every $\nu\in\cM$ there is $c>0$ with $M\ge c\nu$. A universal
semimeasure dominates the uniform measure $\lambda(x)=2^{-|x|}$, so
$M(x)>0$ for all $x$.

\begin{lemma}[root lifting]\label{lem:root}
For $\nu\in\cM$ let $\nu^\uparrow(\eps)=1$ and $\nu^\uparrow(x)=\nu(x)$ for
$x\neq\eps$. Then $\nu^\uparrow\in\cMone$ and $\nu\le\nu^\uparrow$. Hence a
semimeasure dominates every member of $\cM$ if and only if it dominates every
member of $\cMone$, and an lsc semimeasure $M$ is universal if and only if
$M^\uparrow$ is.
\end{lemma}

\begin{proof}
Raising the value at $\eps$ to $1$ keeps the inequality
$\nu(0)+\nu(1)\le\nu(\eps)\le1$, and lower semicomputability is unchanged.
If $M\ge c\nu^\uparrow$ for all $\nu$, then $M\ge c\nu$. If $M\ge c\nu$ for all
$\nu\in\cM$, then for $\nu\in\cMone$ this is the claim. Finally, let $M$ be
lsc. Then $M^\uparrow$ is lsc. If $M$ or $M^\uparrow$ is universal, it
dominates $\lambda$, so $M(\eps)\ge M(0)>0$. Then $M^\uparrow\ge M$ and
$M^\uparrow\le M/M(\eps)$, so the two dominate the same semimeasures.
\end{proof}

All predictions below compare or normalize values $s(x0),s(x1)$ of a
semimeasure at nonempty strings. They never read $s(\eps)$. So
Lemma~\ref{lem:root} lets us state the formal results for $\cMone$ without
loss. The library~\cite{milovanovlean} works with $\cMone$; we say so where it
matters.

\paragraph{Prefix complexity.}
Fix an optimal prefix machine $U$. $K(x)$ and $K(y|x)$ are the prefix
complexity and the conditional prefix complexity, where the program for $y$
receives all of $x$ (so $|x|$ is free). For $t\in\N$, $K(t)$ is the complexity
of a standard binary code of $t$, and $\kappa_t:=2^{-K(t)}$. We use three standard
facts~\cite{livitanyi}:
\begin{align}
&\textstyle\sum_{t}\kappa_t\le1 \quad\text{(Kraft)}, \label{eq:kraft}\\
&\kappa_t\ge c_0(t+1)^{-2}\ \text{for a constant } c_0>0, \label{eq:wlower}\\
&g:\N\to[0,\infty)\ \text{lsc with } \textstyle\sum_n g(n)\le1
\ \Longrightarrow\ g(n)\le C_g\,\kappa_n\ \text{for all } n \quad\text{(coding theorem)}. \label{eq:coding}
\end{align}
The function $t\mapsto \kappa_t$ is lsc. For a computable $F:\Bs\to\Bs$ there is
a constant $b_F$ with $K(F(x)|x)\le b_F$ for all $x$.

\paragraph{Bayes decisions for 0-1 loss.}
For a semimeasure $s$ with $s>0$ and a history $h$, the Bayes action for
binary 0-1 loss compares $s(h0)$ and $s(h1)$. We write $\Lambda_s(h)=1$ if
$s(h1)>s(h0)$, $\Lambda_s(h)=0$ if $s(h1)<s(h0)$, and apply a fixed tie rule
if they are equal. ``Ties toward 1'' means $\Lambda_s(h)=1$ if and only if
$s(h1)\ge s(h0)$. This is the book's action $y^{\Lambda_s}$ of eq.~(3.45) for
0-1 loss; it is scale invariant, so it is the same for $s$ and for its
normalization
\[
q_s(b|h):=\frac{s(hb)}{s(h0)+s(h1)} .
\]
With ties toward 1, $\Lambda_s(h)=1$ if and only if $q_s(1|h)\ge 1/2$.
For a measure $\mu$ the expected cumulative loss (book eq.~(3.46)) is
\[
L_n^{s}:=\E_\mu\sum_{t=1}^n \ind\{\Lambda_s(X_{<t})\neq X_t\}
=\sum_{t=1}^n\sum_{|h|=t-1}\mu(h)\,r_0(h,\Lambda_s(h)),
\]
with the one-step risk $r_0(h,1)=\mu(0|h)$ and $r_0(h,0)=\mu(1|h)$.
We write $L^s_n$ for the book's $L_n^{\Lambda_s}$.

\paragraph{Two inequalities.}
For binary distributions $p=(1-p_1,p_1)$, $q=(1-q_1,q_1)$ with $q_1\in(0,1)$:
\begin{equation}\label{eq:chi2}
\KL(p\Vert q)\ge 2(p_1-q_1)^2 \quad\text{(Pinsker)},\qquad
\frac{p_0^2}{q_0}+\frac{p_1^2}{q_1}=1+\frac{(p_1-q_1)^2}{q_0q_1}\ge1 .
\end{equation}
For probability laws $P,Q$ on $\B^N$ given by conditionals $p,q$,
the chain rule gives
\[
\KL(P\Vert Q)=\sum_{t=1}^N\E_P\,\KL(p(\cdot|X_{<t})\Vert q(\cdot|X_{<t})),
\]
and the hybrid argument (replace the conditionals of $Q$ by those of $P$ one time
step at a time) gives
\begin{equation}\label{eq:hybrid}
\|P-Q\|_1\le\sum_{t=1}^N\E_P\,\|p(\cdot|X_{<t})-q(\cdot|X_{<t})\|_1 .
\end{equation}

%======================================================================
\section{Problem 3.3: two mixtures with unbounded loss difference}\label{sec:33}
%======================================================================

The known part of Problem 3.3 states
$\sum_t\E[\sum_{x_t}(\xi(x_t|x_{<t})-\xi'(x_t|x_{<t}))^2]\le2[\ln w_\mu^{-1}+\ln (w'_\mu)^{-1}]$
and $L_n^{\xi}-L_n^{\xi'}\le O\bigl((L^{\mu}_n)^{1/2}\bigr)$. The predictions of $\xi$
and $\xi'$ merge. The 0-1 decision is a threshold of the prediction at
$1/2$. If $\mu(1|x_{<t})$ approaches $1/2$ at a rate $\delta_t$ with
$\sum_t\delta_t^2<\infty$ and $\sum_t\delta_t=\infty$, the data never reveal
whether the bias is $+\delta_t$ or $-\delta_t$ (a Kakutani-type
indistinguishability), and two priors that favor different signs keep making
different decisions. Each different decision costs $2\delta_t$ in expectation,
and the sum diverges. The same mechanism appears in the book's
Section~3.6 and Problem~3.6~\cite[pp.~98, 120]{hutter2005uai} and in
Hutter~\cite[p.~988]{hutter2003jmlr}, where a class $\{\mu_1,\mu_2\}$ with
$\mu(1|x_{<t})=\frac12\pm\varepsilon_t$ shows that the regret of \emph{one}
mixture against the informed predictor can diverge. Problem 3.3 compares
\emph{two} mixtures; our first example is in the same family.

\subsection{A two-measure example}

\begin{theorem}[two measures]\label{thm:warmup}
Let $\delta_t=1/(32t)$. Let $\mu$ and $\nu$ be the product measures with
$\mu(1|x_{<t})=\frac12+\delta_t$ and $\nu(1|x_{<t})=\frac12-\delta_t$. Let
$\zeta_-=\frac14\mu+\frac34\nu$ and $\zeta_+=\frac34\mu+\frac14\nu$, two Bayes
mixtures over $\cM=\{\mu,\nu\}$ with positive weights. With 0-1 loss and ties
toward $0$, for every $n\ge1$,
\[
\frac{H_n}{32}\le L_n^{\zeta_-}-L_n^{\zeta_+}\le\frac{H_n}{16},
\qquad H_n=\sum_{t=1}^n\frac1t .
\]
In particular the difference is unbounded. Both measures are computable with
rational conditionals. (Formal version: the lower bound and unboundedness; the
upper bound is a paper step.)
\end{theorem}

\begin{proof}
Fix $t$ and a history $h$ of length $t-1$; put $a=\mu(h)>0$, $b=\nu(h)>0$,
$\delta=\delta_t$. For a mixture $\zeta=w_\mu\mu+w_\nu\nu$,
$\zeta(h1)-\zeta(h0)=2\delta(w_\mu a-w_\nu b)$. With ties toward 0,
$\Lambda_{\zeta_-}(h)=1$ if and only if $a>3b$, and $\Lambda_{\zeta_+}(h)=1$ if and
only if $3a>b$. So $\Lambda_{\zeta_-}(h)=1$ implies $\Lambda_{\zeta_+}(h)=1$. The
risks under $\mu$ are $r(h,1)=\frac12-\delta$ and $r(h,0)=\frac12+\delta$, so
$g(h):=r(h,\Lambda_{\zeta_-}(h))-r(h,\Lambda_{\zeta_+}(h))$ is $2\delta$ if
$b/3<a\le 3b$ (the decisions differ) and $0$ otherwise. We claim
\begin{equation}\label{eq:warmup-pointwise}
a\,g(h)\ge2\delta\Bigl(a-\frac{4(b-a)^2}{a}\Bigr).
\end{equation}
If the decisions differ, the left side is $2\delta a$. If they agree, then
$a>3b$ or $b\ge3a$; in both cases $|b-a|\ge a/2$, and the right side is
$\le0$.

Let $R(h)=\nu(h)/\mu(h)$. Summing~\eqref{eq:warmup-pointwise} over $|h|=t-1$
(its left side is $\mu(h)g(h)$) gives
\[
\sum_{|h|=t-1}\mu(h)g(h)\ge2\delta_t\bigl(1-4\,\E_\mu[(R-1)^2]\bigr).
\]
Since $\E_\mu R=1$ and the bits are independent,
\[
\E_\mu R^2=\prod_{s<t}\biggl(\sum_{c\in\B}\frac{\nu_s(c)^2}{\mu_s(c)}\biggr)
=\prod_{s<t}\Bigl(1+\frac{16\delta_s^2}{1-4\delta_s^2}\Bigr)
\le\exp\biggl(\sum_{s\ge1}\frac{16}{1020\,s^2}\biggr)<1.03<\frac98,
\]
using~\eqref{eq:chi2} for each factor and $1-4\delta_s^2\ge 1020/1024$.
So $\E_\mu[(R-1)^2]=\E_\mu R^2-1\le\frac18$, and the expected one-step
difference at time $t$ is at least $\delta_t$. Summing over $t\le n$ gives the
lower bound. Each one-step difference is at most $2\delta_t$, which gives the
upper bound.
\end{proof}

The class here is a class of computable
probability measures, which is the setting of the book's Chapter~3
(cf.\ Problem 3.1). The rest of this section strengthens the example to
universal mixtures. Section~\ref{sec:33measures} then returns to classes of
computable probability measures, with an arbitrary positive base prior.

\subsection{A finite experiment that survives any small background}

A universal mixture contains every computable environment, so it cannot be
kept away from the decision boundary by choosing the class. The core of the
construction is a finite experiment whose decision gap survives any added
background semimeasure of relative weight at most $2^{-20}$.

\begin{theorem}[robust finite experiment]\label{thm:robust}
Let $k\ge1$, $N=k^2$ and $\gamma=1/(32k)$. Let $\mu^+_N,\mu^-_N$ be the measures
with independent bits, $\mu^+_N(1|h)=\frac12+\gamma$ and $\mu^-_N(1|h)=\frac12-\gamma$
for $|h|<N$, and fair bits after time $N$. Let
$\pi_-=\frac14\mu^+_N+\frac34\mu^-_N$ and $\pi_+=\frac34\mu^+_N+\frac14\mu^-_N$.
Let $A$ be any semimeasure with $A(\eps)=1$, let $0\le\eta\le2^{-20}$, and let
$s_\pm=(1-\eta)\pi_\pm+\eta A$. With 0-1 loss, ties toward 1, and losses under
$\mu^+_N$,
\[
\frac{k}{32}\le L_N^{s_-}-L_N^{s_+}\le\frac{k}{16},
\]
and every one-step expected loss difference is nonnegative.
\end{theorem}

\begin{proof}
Write $\mu^\pm=\mu^\pm_N$ and $R_m=\mu^-(X_{1:m})/\mu^+(X_{1:m})$.

\emph{Step 1 (most histories have a margin).} For $m<N$, $\E_{\mu^+} R_m=1$ and,
as in the proof of Theorem~\ref{thm:warmup},
\[
\E_{\mu^+} R_m^2=\Bigl(1+\frac{16\gamma^2}{1-4\gamma^2}\Bigr)^m
\le\exp\Bigl(\frac{16N\gamma^2}{1-4\gamma^2}\Bigr)\le e^{4/255}<\frac{255}{251},
\]
because $N\gamma^2=1/1024$ and $1-4\gamma^2\ge255/256$. By Chebyshev's
inequality,
$\Pr_{\mu^+}\{R_m\notin[\frac12,2]\}\le\Pr_{\mu^+}\{|R_m-1|\ge\frac12\}\le4(\E_{\mu^+} R_m^2-1)\le\frac{16}{251}$.
Call a history $h$ with $|h|<N$ \emph{good} if $R(h)=\mu^-(h)/\mu^+(h)\in[\frac12,2]$.
A direct computation gives
\[
q_{\pi_-}(1|h)=\frac12+\gamma\,\frac{1-3R(h)}{1+3R(h)},\qquad
q_{\pi_+}(1|h)=\frac12+\gamma\,\frac{3-R(h)}{3+R(h)} .
\]
Both fractions are decreasing in $R$, so on good histories
$q_{\pi_-}(1|h)\le\frac12-\frac\gamma5$ and $q_{\pi_+}(1|h)\ge\frac12+\frac\gamma5$.
The expected number of good rounds among $t=1,\dots,N$ is at least
$(1-\frac{16}{251})N>\frac34N$.

\emph{Step 2 (the background moves few decisions).}
Let $\bar s$ be the probability measure with conditionals $q_s$. Since
$s(h0)+s(h1)\le s(h)$, the product of the conditionals along $x$ gives
$\bar s(x)\ge s(x)/s(\eps)=s(x)$, and $s(\eps)=1$. So
$\bar s_\pm\ge s_\pm\ge(1-\eta)\pi_\pm$ on strings of length $\le N$, and
$\KL(\pi_\pm\Vert\bar s_\pm)\le\ln\frac1{1-\eta}$ on $\B^N$. The chain rule
and Pinsker's inequality~\eqref{eq:chi2} give
\[
\sum_{t=1}^N\E_{\pi_\pm}\bigl(q_{s_\pm}(1|X_{<t})-q_{\pi_\pm}(1|X_{<t})\bigr)^2
\le\frac12\ln\frac1{1-\eta}.
\]
Since $\pi_\pm\ge\frac14\mu^+$, the same sum under $\mu^+$ is at most
$2\ln\frac1{1-\eta}$. By Markov's inequality, the expected number of rounds in
which $|q_{s_\pm}-q_{\pi_\pm}|\ge\gamma/5$ for at least one of the two signs is
at most
\[
\frac{100}{\gamma^2}\ln\frac1{1-\eta}\le\frac{200\eta}{\gamma^2}
=200\eta\cdot1024N\le\frac{25}{128}N .
\]
On a good round where neither forecast moved by $\gamma/5$, we have
$q_{s_-}(1|h)<\frac12<q_{s_+}(1|h)$, so $\Lambda_{s_-}(h)=0$ and
$\Lambda_{s_+}(h)=1$. The expected number of such rounds is at least
$(\frac34-\frac{25}{128})N=\frac{71}{128}N>\nobreak\frac N2$.

\emph{Step 3 (no negative steps).} For $|h|<N$,
\[
[s_+(h1)-s_+(h0)]-[s_-(h1)-s_-(h0)]=(1-\eta)\gamma\bigl(\mu^+(h)+\mu^-(h)\bigr)>0;
\]
the term $\eta A$ cancels. So $\Lambda_{s_-}(h)=1$ implies
$s_+(h1)>s_+(h0)$ and $\Lambda_{s_+}(h)=1$. The decisions differ only as
$(\Lambda_{s_-},\Lambda_{s_+})=(0,1)$, and such a round adds exactly $2\gamma$
to the expected loss difference under $\mu^+$; all other rounds add $0$. With
Step 2, $L_N^{s_-}-L_N^{s_+}\ge2\gamma\cdot\frac N2=\frac{k}{32}$, and
$L_N^{s_-}-L_N^{s_+}\le2\gamma N=\frac k{16}$.
\end{proof}

\begin{remark}
The proof of Theorem~\ref{thm:robust} does not use the tie rule: Step 3 gives
$s_+(h1)>s_+(h0)$ strictly whenever $s_-(h1)\ge s_-(h0)$, and in Step 2 a tie
of $s_\pm$ also needs a forecast change of at least $\gamma/5$. The formal
statement fixes ties toward 1. The bound holds for every background $A$, with
no rate of convergence for $A$. Theorems~\ref{thm:aggregate},
\ref{thm:mixture}, \ref{thm:complexity}, \ref{thm:exact} and~\ref{thm:measures} take their decisions from
Theorem~\ref{thm:robust} (outside the experiments both actions have risk
$\frac12$), so they also hold for every fixed tie rule; their formal
statements fix ties toward 1.
\end{remark}

\subsection{Aggregation over a universal background}\label{sec:33aggregate}

Next we put infinitely many experiments behind a selector and make two
semimeasures copy the background's mass on each selector. Then the relative
weight of the background is exactly $\eta$ in every branch, however small the
branch's mass is.

\paragraph{The source.}
For $j\ge0$ let $z_j=1^j0$, $p_j=2^{-j-1}$, $N_j=16^j$ and $\gamma_j=1/(32\cdot4^j)$,
and let $\mu^\pm_{N_j}$ be the measures of
Theorem~\ref{thm:robust} with $k=4^j$. The source $\mu$ produces fair bits
until the first $0$; if this $0$ follows $j$ ones, it produces $N_j$ bits with
$\mu(1|\cdot)=\frac12+\gamma_j$, and then fair bits forever. Every string is
either $1^k$ or $z_jh$ for unique $j,h$, and
\[
\mu(1^k)=2^{-k},\qquad \mu(z_jh)=p_j\,\mu^+_{N_j}(h).
\]
All conditionals of $\mu$ are rational and computable from the history.

\paragraph{The copies.}
For a semimeasure $M$ define
\[
M^\pm(1^k)=M(1^k),\qquad M^\pm(z_jh)=M(z_j)\,\mu^\pm_{N_j}(h).
\]
$M^+$ is a semimeasure: at $1^k$,
$M^+(1^k0)+M^+(1^{k+1})=M(z_k)+M(1^{k+1})\le M(1^k)$, and inside a branch the
equality comes from $\mu^+_{N_j}$. The same holds for $M^-$. Let $\eta=2^{-20}$ and
\[
\xi_-=(1-\eta)\bigl(\tfrac14M^{+}+\tfrac34M^{-}\bigr)+\eta M,\qquad
\xi_+=(1-\eta)\bigl(\tfrac34M^{+}+\tfrac14M^{-}\bigr)+\eta M .
\]

\begin{theorem}[aggregation]\label{thm:aggregate}
Let $M$ be a semimeasure with $M(x)>0$ for all $x$. Then $\xi_\pm$ are
semimeasures with $\xi_\pm\ge\eta M$, and with 0-1 loss, ties toward 1 and
losses under the source $\mu$, for every $n\ge2$,
\[
\frac{n^{1/4}}{256}\le L_n^{\xi_-}-L_n^{\xi_+}\le\frac{n^{1/4}}8 .
\]
No computability of $M$ is needed.
\end{theorem}

\begin{proof}
The semimeasure property and $\xi_\pm\ge\eta M$ are clear. Consider a round $t$
with history $g$, $|g|=t-1$.
If $g=1^k$, or $g=z_jh$ with $|h|\ge N_j$, then $\mu(1|g)=\frac12$, so both
predictors have risk $\frac12$ and the round adds nothing.
If $g=z_jh$ with $|h|<N_j$, then, since $M^\pm(z_j)=M(z_j)>0$,
\[
\frac{\xi_\pm(z_jhb)}{M(z_j)}=(1-\eta)\pi^\pm_j(hb)+\eta A_j(hb),\qquad
A_j(y):=\frac{M(z_jy)}{M(z_j)},
\]
where $\pi^\pm_j$ are the cores of Theorem~\ref{thm:robust} with $k=4^j$, and
$A_j$ is a semimeasure with $A_j(\eps)=1$. Bayes actions are scale invariant,
so $\Lambda_{\xi_\pm}(z_jh)=\Lambda_{s^\pm_j}(h)$ with $s^\pm_j$ as in
Theorem~\ref{thm:robust} for $A=A_j$. Also $\mu(z_jh)=p_j\mu^+_{N_j}(h)$ and
$\mu(\cdot|z_jh)=\mu^+_{N_j}(\cdot|h)$. Let $D_j(m)$ be the loss difference of
Theorem~\ref{thm:robust} (with $k=4^j$, $A=A_j$) over its first $m$ rounds.
By the theorem, $D_j$ is nondecreasing, $D_j(N_j)\ge4^j/32$, and
$D_j(m)\le2\gamma_jm$. A round $t\le n$ has history $z_jh$ with $|h|<N_j$ for
exactly ${m_j(n):=\min\{N_j,(n-j-1)_+\}}$ values of $t$, so
\[
\Delta_n:=L_n^{\xi_-}-L_n^{\xi_+}=\sum_{j\ge0}p_jD_j(m_j(n)).
\]
\emph{Lower bound.} Let $J$ be the largest integer with $2\cdot16^J\le n$
(so $J\ge0$). For $j\le J$, $j+1+16^j\le2\cdot16^j\le n$, so branch $j$ is
complete and contributes at least $p_j\,4^j/32=2^j/64$. Hence
$\Delta_n\ge(2^{J+1}-1)/64\ge2^J/64$. From $n<2\cdot16^{J+1}$ we get
$2^J>n^{1/4}/(2\cdot2^{1/4})$, so $\Delta_n\ge n^{1/4}/(128\cdot2^{1/4})\ge n^{1/4}/256$.

\emph{Upper bound.}
$\Delta_n\le\sum_jp_j2\gamma_j\min(16^j,n)=\frac1{32}\sum_j8^{-j}\min(16^j,n)$.
With $J_0=\lfloor\log_{16}n\rfloor$, the terms $j\le J_0$ sum to at most
$2^{J_0+1}\le2n^{1/4}$ and the terms $j>J_0$ to at most
$\frac87n\,8^{-J_0-1}\le\frac87n^{1/4}$. So
$\Delta_n\le\frac{22}{7\cdot32}n^{1/4}<\frac18n^{1/4}$.
\end{proof}

Each path selects one finite experiment, so for this pair the
\emph{realized} difference of error counts is eventually constant on every
path. The expected difference diverges because the experiment lengths have a
heavy tail. Problem 3.3 asks about expected losses.

\subsection{Two Bayes mixtures over all lower semicomputable semimeasures}

\begin{theorem}[main theorem for Problem 3.3]\label{thm:mixture}
Let $\nu_0,\nu_1,\dots$ be a uniformly lower semicomputable enumeration, with
repetitions, of all lower semicomputable semimeasures with $\nu(\eps)=1$.
There are weights $w^-,w^+$ with $w^\pm_i>0$, $\sum_iw^\pm_i=1$ and
$\frac13\le w^-_i/w^+_i\le3$ for all $i$, such that both mixtures
$\xi_\pm=\sum_iw^\pm_i\nu_i$ are lower semicomputable semimeasures, and an
index $i_0$ such that $\nu_{i_0}=\mu$ is a computable probability measure with
rational conditionals. For binary 0-1 loss, with ties broken toward 1,
\[
L_n^{\xi_-}-L_n^{\xi_+}\ge\frac{n^{1/4}}{256}\quad\text{for all }n\ge2,
\]
and the difference tends to $+\infty$. Both mixtures are universal (paper step).
\end{theorem}

\begin{proof}
Let $M=\sum_i2^{-(i+1)}\nu_i$. It is an lsc semimeasure with $M(\eps)=1$; it
dominates every $\nu_i$, hence (Lemma~\ref{lem:root}) every member of $\cM$;
and $M>0$ because $\lambda$ is some $\nu_i$. Build $\mu,M^+,M^-$ from this $M$ as
above. The order of the construction is: the enumeration, then $M$, then $M^+$
and $M^-$, then the weights; $M$ does not depend on the weights.

$M^\pm$ are lower semicomputable: $M^\pm(x)=c_\pm(x)\,M(\sigma(x))$, where
$\sigma(x)$ is the prefix $1^k$ or $z_j$ of $x$ and $c_\pm(x)\in\mathbb Q$ is
computable ($c_\pm(1^k)=1$, $c_\pm(z_jh)=\mu^\pm_{N_j}(h)$). Also $M^\pm(\eps)=M(\eps)=1$.
So $M^\pm\in\cMone$. Also $M^+\neq M^-$, because
$M^+(z_01)=M(z_0)(\frac12+\gamma_0)\neq M(z_0)(\frac12-\gamma_0)=M^-(z_01)$ and
$M(z_0)>0$. So there are indices $i_+\neq i_-$ with $\nu_{i_+}=M^+$, $\nu_{i_-}=M^-$.
Put
\begin{gather*}
w^-_i=\eta2^{-(i+1)}+(1-\eta)\bigl(\tfrac14\ind\{i=i_+\}+\tfrac34\ind\{i=i_-\}\bigr),\\
w^+_i=\eta2^{-(i+1)}+(1-\eta)\bigl(\tfrac34\ind\{i=i_+\}+\tfrac14\ind\{i=i_-\}\bigr).
\end{gather*}
All weights are positive and each family sums to $\eta+(1-\eta)=1$. Each of the
three summands of $w^-_i$ is between $\frac13$ and $3$ times the matching summand
of $w^+_i$, so the same holds for $w^-_i/w^+_i$. By construction
$\sum_iw^\pm_i\nu_i$ equals $\xi_\pm$ of Theorem~\ref{thm:aggregate}, so it is
an lsc semimeasure, it dominates $\eta M$ and is therefore universal, and the
loss bound holds. The source $\mu$ is a computable measure, hence in $\cMone$,
so $\mu=\nu_{i_0}$ for some $i_0$, and $w^\pm_{i_0}>0$.
\end{proof}

\begin{remark}[the book's class $\cM_U$]\label{rem:MU}
The proof uses only three facts about the enumeration: it is uniformly lsc, it
contains $M^+$, $M^-$, $\mu$ and $\lambda$, and Theorem~\ref{thm:aggregate}
accepts any positive background. So the same proof works for any uniformly lsc
enumeration of the book's class $\cM_U$ of all lsc semimeasures with
$\nu(\eps)\le1$: then $\xi_\pm$ are Bayes mixtures over $\cM_U$ in the sense
of eq.~(3.5) of the book, with positive weights, and $\mu\in\cM_U$. The formal
proof uses the enumeration of $\cMone$ in~\cite{milovanovlean}.
\end{remark}

\begin{remark}[what is not claimed]
The weights of Theorem~\ref{thm:mixture} are computable: they are rational,
and they differ from the computable sequence $\eta2^{-(i+1)}$ only at the two
fixed indices $i_+$ and $i_-$ (a paper observation; the formal statement does not
include it). Theorem~\ref{thm:mixture} does not say that the weights are of
the form $2^{-K(\nu)}$, or that $\xi_\pm$ are Solomonoff priors $M_U$ of
universal monotone machines. Weights within constant factors of $2^{-K(i)}$
are the subject of Theorem~\ref{thm:complexity}, and exact weights
$2^{-K_C(i)}$ that of Theorem~\ref{thm:exact}. Theorem~\ref{thm:mixture}
also says nothing about the specific product source of
Theorem~\ref{thm:warmup} inside a universal mixture.
\end{remark}

\subsection{Complexity weights $2^{-K(i)}$}

The book's universal prior uses weights $w_\nu=2^{-K(\nu)}$, with
$K(\nu_i):=K(i)$ for the enumeration index~\cite[pp.~48, 67]{hutter2005uai}.

\begin{theorem}[complexity priors]\label{thm:complexity}
Fix an optimal prefix machine and the enumeration $(\nu_i)$ of
Theorem~\ref{thm:mixture}. There are weights $w^-,w^+$ and constants
$0<c\le C<\infty$ (depending on the machine) with
$c\,2^{-K(i)}\le w^\pm_i\le C\,2^{-K(i)}$, where $K(i)$ is the prefix
complexity of the index $i$, with $\sum_iw^\pm_i\le1$ and
$w^-_i/w^+_i\in[\frac13,3]$, such that both mixtures
$\xi^K_\pm=\sum_iw^\pm_i\nu_i$ are lower semicomputable semimeasures, and an index
$i_0$ with $\nu_{i_0}=\mu$ a computable probability measure. For binary 0-1 loss,
ties toward 1,
\[
L_n^{\xi^K_-}-L_n^{\xi^K_+}\ge\frac{n^{1/4}}{256}\quad\text{for all }n\ge2,
\]
and the difference tends to $+\infty$. Both mixtures are universal (paper step).
\end{theorem}

\begin{proof}
Let $M_K=\sum_i2^{-K(i)}\nu_i$. It is an lsc semimeasure (since $2^{-K(i)}$
is lsc and $\sum_i2^{-K(i)}\le1$ by~\eqref{eq:kraft}), it is positive, and it
has $M_K(\eps)<1$ in general. Build $M_K^\pm$ from $M_K$ as $M^\pm$ was built from $M$. Their roots are
$M_K(\eps)$, so we use the lifts $(M_K^\pm)^\uparrow\in\cMone$ of
Lemma~\ref{lem:root}. As in the proof of Theorem~\ref{thm:mixture},
$(M_K^+)^\uparrow(z_01)\neq(M_K^-)^\uparrow(z_01)$ because $M_K(z_0)>0$, so they have
indices $i_+\neq i_-$. Put
\[
w^\pm_i=\eta2^{-K(i)}+(1-\eta)\bigl(a_\pm\ind\{i=i_+\}+b_\pm\ind\{i=i_-\}\bigr),
\quad (a_-,b_-)=(\tfrac14,\tfrac34),\ (a_+,b_+)=(\tfrac34,\tfrac14).
\]
Then $w^\pm_i\ge\eta2^{-K(i)}$, so $c=\eta$ works; $w^\pm_i=\eta2^{-K(i)}$ for
$i\notin\{i_+,i_-\}$ and $w^\pm_i\le\eta2^{-K(i)}+1$ for $i\in\{i_+,i_-\}$, so
$C=\eta+\max(2^{K(i_+)},2^{K(i_-)})$ works; $\sum_iw^\pm_i\le\eta+(1-\eta)=1$; and the
ratio bound holds term by term. The mixture is
$\xi^K_\pm=\sum_iw^\pm_i\nu_i=\eta M_K+(1-\eta)(a_\pm(M_K^+)^\uparrow+b_\pm(M_K^-)^\uparrow)$. It is
an lsc semimeasure: its root value is $\eta M_K(\eps)+(1-\eta)\le1$, and
$\xi^K_\pm(0)+\xi^K_\pm(1)\le\eta M_K(\eps)+(1-\eta)M_K(\eps)\le\eta M_K(\eps)+(1-\eta)=\xi^K_\pm(\eps)$
because $M_K(\eps)\le1$; at other strings the inequality holds termwise. It
agrees on every nonempty string with the $\xi_\pm$ of
Theorem~\ref{thm:aggregate} built from the positive background $M_K$. The
losses read only nonempty strings, so Theorem~\ref{thm:aggregate} gives the
bound. Universality follows from $\xi^K_\pm\ge\eta2^{-K(i)}\nu_i$ for every $i$
and Lemma~\ref{lem:root}.
\end{proof}

The weights here are per index. We do not claim
$w_\nu\le C\,2^{-K(\nu)}$ for the total weight of one semimeasure that occurs
at several indices, or $\sum_iw_i=1$.

In the book, the universal prior has the exact weights $2^{-K(i)}$, and $K$
is defined relative to a choice of optimal prefix machine. The next theorem
keeps these exact weights and changes only the reference machine. For a prefix
machine $C$ and an enumeration $(\tau_i)$, let
\[
\xi_C:=\sum_i2^{-K_C(i)}\,\tau_i ,
\]
where $K_C(i)$ is the prefix complexity of the index $i$ relative to $C$.

\begin{theorem}[exact complexity weights]\label{thm:exact}
There are a uniformly lower semicomputable enumeration
$\tau_0,\tau_1,\dots$, with repetitions, of all lower semicomputable
semimeasures with $\nu(\eps)=1$, and two optimal prefix machines $C_-$, $C_+$,
with these properties. The enumeration is a computable renumbering of the
enumeration $(\nu_i)$ of Theorem~\ref{thm:mixture} in both directions: there
are computable $f,g$ with $\tau_{f(i)}=\nu_i$ and $\nu_{g(i)}=\tau_i$ for all
$i$. The mixtures $\xi_{C_-}$ and $\xi_{C_+}$ are universal lower
semicomputable semimeasures, and
$\frac14\le2^{-K_{C_-}(i)}/2^{-K_{C_+}(i)}\le4$ for all $i$. The source $\mu$
of Section~\ref{sec:33aggregate} is some $\tau_j$, and for binary 0-1 loss,
with ties broken toward 1,
\[
L_n^{\xi_{C_-}}-L_n^{\xi_{C_+}}\ge\frac{n^{1/4}}{256}\quad\text{for all }n\ge2,
\]
and the difference tends to $+\infty$. (Formal version: for the library
enumeration as $(\nu_i)$, with the weaker bound $n^{1/4}/512$ for $n\ge3$.)
\end{theorem}

\begin{proof}
Fix an optimal prefix machine $U_0$. The order of the construction is: the
background $M_0$, then the enumeration $\tau$, then the machines $C_\pm$.

\emph{The background.} Let
$R=\sum_i2^{-30-K_{U_0}(i+4)}\nu_i$ and $M_0=(2^{21}-2)R$. Since
$2^{-K_{U_0}(i)}$ is lsc and $\sum_i2^{-K_{U_0}(i)}\le1$
by~\eqref{eq:kraft}, $M_0$ is an lsc semimeasure with
$M_0(\eps)\le2^{21-30}<1$; it is positive because $\lambda$ is some $\nu_i$.

\emph{The enumeration.} Build the copies $M_0^\pm$ of
Section~\ref{sec:33aggregate} and let $\theta_\pm=(M_0^\pm)^\uparrow\in\cMone$
(Lemma~\ref{lem:root}). Put
\[
\tau_0=\tau_1=\theta_+,\qquad \tau_2=\tau_3=\theta_-,\qquad \tau_{i+4}=\nu_i .
\]
This family is uniformly lsc and contains every $\nu_i$. Take $f(i)=i+4$, and
let $g$ send $0,1$ and $2,3$ to fixed indices of $\theta_+$ and $\theta_-$ in
$(\nu_i)$, and $i+4$ to $i$. The source is $\mu=\nu_{i_0}=\tau_{i_0+4}$.

\emph{The machines.} On every condition $y$, $C_-$ maps the four programs
$0000$, $0001$, $01$, $001$ to the indices $0,1,2,3$, and $C_+$ maps
$01$, $001$, $0000$, $0001$ to them. On the programs $1^{30}p$ both machines
compute $U_0(p|y)$. They have no other programs. The four short programs form a prefix-free set, they start with $0$, and the
programs of $U_0$ form a prefix-free set, so the domain is prefix-free, and
$K_{C_\pm}(x|y)\le K_{U_0}(x|y)+30$, so $C_\pm$ are optimal. The only
programs of length less than $30$ are the four short ones, and a program
$1^{30}p$ outputs $i$ only if $U_0(p)=i$. So $K_{C_-}(i)$ is $4,4,2,3$ and
$K_{C_+}(i)$ is $2,3,4,4$ for $i=0,1,2,3$, and
$K_{C_\pm}(i)=30+K_{U_0}(i)$ for $i\ge4$. The weight ratios are
$\frac14,\frac12,4,2$ at $i=0,1,2,3$ and $1$ for $i\ge4$. With these weights
\[
\xi_{C_-}=\tfrac18\theta_++\tfrac38\theta_-+R,\qquad
\xi_{C_+}=\tfrac38\theta_++\tfrac18\theta_-+R .
\]

\emph{The mixtures.} Each $\tau_i$ is a semimeasure and
$\sum_i2^{-K_{C_\pm}(i)}\le1$, so $\xi_{C_\pm}$ are semimeasures; they are
lsc because $i\mapsto2^{-K_{C_\pm}(i)}$ is lsc and $\tau$ is uniformly lsc.
They are universal, since $\xi_{C_\pm}\ge2^{-K_{C_\pm}(i)}\tau_i$ for every
$i$ and the $\tau_i$ include all of $\cMone$ (Lemma~\ref{lem:root}).

\emph{The loss.} Let $\eta=2^{-20}$, so $(2^{21}-2)/2^{20}=2(1-\eta)$ and
$\eta M_0=2(1-\eta)R$. On nonempty strings $\theta_\pm=M_0^\pm$, so there
\[
2(1-\eta)\,\xi_{C_-}=(1-\eta)\bigl(\tfrac14M_0^++\tfrac34M_0^-\bigr)+\eta M_0,
\]
which is the $\xi_-$ of Theorem~\ref{thm:aggregate} for the positive
background $M=M_0$; in the same way $2(1-\eta)\,\xi_{C_+}$ is its $\xi_+$.
Bayes actions read only nonempty strings and do not change under a positive
factor. So the losses are those of Theorem~\ref{thm:aggregate}, which gives
the bound.
\end{proof}

The two copies of $\theta_+$ and of $\theta_-$ make the total weights
$\frac18$ and $\frac38$ sums of powers of two with prefix-free codes, so the
two foreground weights have the ratio $1:3$ of Theorem~\ref{thm:aggregate}.
Any tag $1^L$ with $L\ge21$ works; the formal proof uses $L=30$. The theorem
holds in the index convention $K(\nu_i):=K(i)$. It does not say that every
pair of optimal machines gives a gap (with $C_-=C_+$ the difference is $0$),
and it makes no claim for the enumeration $(\nu_i)$ itself: $\tau$ adds four
entries in front of it. Semantic weights $2^{-K(\nu)}$ are the subject of
Remark~\ref{rem:exactK}. Solomonoff's mixtures $M_U$ of universal monotone
machines are not covered.

\subsection{Two Bayes mixtures over a class of probability measures}\label{sec:33measures}

In the book's Section~3.2.2, $\cM$ is a countable class of probability
measures, $\mu\in\cM$, and $\xi=\sum_{\nu\in\cM}w_\nu\nu$ with $w_\nu>0$ and
$\sum_\nu w_\nu=1$~\cite[p.~70, eq.~(3.5)]{hutter2005uai}. Problem~3.3 is
stated in this setting. The classes of Theorems~\ref{thm:mixture},
\ref{thm:complexity} and~\ref{thm:exact} contain semimeasures that are not measures. This
subsection gives the answer for classes of probability measures, with priors
$w^\pm\ge2^{-20}a$ built from any positive base prior $a$ on the class.

\paragraph{The class.}
For a function $r:\Bs\to\mathbb Q\cap[0,1]$ let $\nu_r$ be the probability
measure with $\nu_r(1|x)=r(x)$: $\nu_r(\eps)=1$, $\nu_r(x1)=\nu_r(x)\,r(x)$
and $\nu_r(x0)=\nu_r(x)(1-r(x))$. Let $\cCQ$ be the family $(\nu_r)_r$, where
$r$ ranges over the total computable functions $\Bs\to\mathbb Q\cap[0,1]$. The
family is indexed by $r$, so one measure can occur at several indices (two
functions $r$ that differ only at $\nu_r$-null histories give the same
measure). A prior on $\cCQ$ is a function $r\mapsto a_r$ on the indices, and
$\cCQ$ is countable. The source $\mu$ of Section~\ref{sec:33aggregate} has total computable
rational conditionals, so $\mu=\nu_{r_0}$ for some index $r_0$; the fair
coin is $\lambda=\nu_{1/2}$.

For an index $r$ let $r^\pm$ copy $r$ on the selector and run the branch
experiment of the source and then fair bits:
\[
r^\pm(1^k)=r(1^k),\qquad
r^\pm(z_jh)=\begin{cases}\frac12\pm\gamma_j,&|h|<N_j,\\[2pt] \frac12,&|h|\ge N_j.\end{cases}
\]
The functions $r^\pm$ are again total, computable and rational, and
$\nu_{r^\pm}(1^k)=\nu_r(1^k)$, $\nu_{r^\pm}(z_jh)=\nu_r(z_j)\,\mu^\pm_{N_j}(h)$. So
$\nu_{r^\pm}=(\nu_r)^\pm$, the copies of Section~\ref{sec:33aggregate}.

\begin{theorem}[mixtures over measures]\label{thm:measures}
There is an index $r_0$ such that $\nu_{r_0}=\mu$ is the source of
Section~\ref{sec:33aggregate}, with the following property. For every prior $a$ on $\cCQ$ with
$a_r>0$ for all $r$ and $\sum_ra_r=1$, there are priors $w^-,w^+$ on $\cCQ$
with
\[
w^\pm_r>0,\qquad \sum_rw^\pm_r=1,\qquad w^\pm_r\ge2^{-20}a_r,\qquad
\tfrac13\le w^-_r/w^+_r\le3\quad\text{for all }r,
\]
such that for the Bayes mixtures $\xi_\pm=\sum_rw^\pm_r\nu_r$, binary 0-1 loss
and ties toward 1,
\[
L_n^{\xi_-}-L_n^{\xi_+}\ge\frac{n^{1/4}}{256}\quad\text{for all }n\ge2,
\]
and the difference tends to $+\infty$. The source $\mu$ does not depend on
$a$; the priors $w^\pm$ depend on $a$. We make no claim that $w^\pm$ or $\xi_\pm$
are computable. The proof also gives the upper bound $n^{1/4}/8$. (Formal
version: as stated; the upper bound is a paper step.)
\end{theorem}

\begin{proof}
Let $M_a=\sum_ra_r\nu_r$. It is a probability measure, and
$M_a\ge a_{1/2}\lambda>0$. For an index $r'$ put
\[
\tilde a^\pm_{r'}=\sum_{r:\,r^\pm=r'}a_r ,
\]
the prior mass of the fiber of $r'$ under $r\mapsto r^\pm$ (zero if the fiber is
empty). The fibers can have more than one element, and the proof does not need
them to be singletons. Summing over fibers gives
$\sum_{r'}\tilde a^\pm_{r'}=\sum_ra_r=1$ and
\[
\sum_{r'}\tilde a^\pm_{r'}\nu_{r'}=\sum_ra_r\,\nu_{r^\pm}=\sum_ra_r(\nu_r)^\pm=M_a^\pm ,
\]
because $M\mapsto M^\pm$ is linear: $M^\pm(x)=c_\pm(x)\,M(\sigma(x))$ with
$\sigma$ and $c_\pm$ as in the proof of Theorem~\ref{thm:mixture}. With
$\eta=2^{-20}$ put
\begin{gather*}
w^-_r=\eta a_r+(1-\eta)\bigl(\tfrac14\tilde a^+_r+\tfrac34\tilde a^-_r\bigr),\\
w^+_r=\eta a_r+(1-\eta)\bigl(\tfrac34\tilde a^+_r+\tfrac14\tilde a^-_r\bigr).
\end{gather*}
Then $w^\pm_r\ge\eta a_r>0$, each family sums to $\eta+(1-\eta)=1$, and the
ratio bound holds term by term. By the display above,
\[
\xi_-=\eta M_a+(1-\eta)\bigl(\tfrac14M_a^++\tfrac34M_a^-\bigr),\qquad
\xi_+=\eta M_a+(1-\eta)\bigl(\tfrac34M_a^++\tfrac14M_a^-\bigr),
\]
which are the $\xi_\pm$ of Theorem~\ref{thm:aggregate} for the positive
background $M=M_a$. Theorem~\ref{thm:aggregate} gives the bound and the upper
bound $n^{1/4}/8$. Finally $w^\pm_{r_0}\ge\eta a_{r_0}>0$.
\end{proof}

The construction transforms every component of the given prior. The copies
keep the selector masses of each component, so the background has relative
weight exactly $\eta$ in every branch. This matters because, with two fixed
foreground models, a small global weight of the background could
become a large weight after conditioning on a rare selector.

\begin{remark}[indices and distinct measures]\label{rem:indices}
The theorem also holds with $\cCQ$ read as a set of distinct measures (paper
step; not formally verified). A positive prior on the distinct measures lifts to a positive prior on
the indices: split the weight of each measure over its countably many indices
with positive parts. Theorem~\ref{thm:measures} then gives $w^\pm$ on indices.
Summing $w^\pm$ over the indices of each measure gives priors on the distinct
measures; positivity, $\sum=1$, the lower bound $2^{-20}a$ and the ratio bound
hold after the summation, and the mixtures $\xi_\pm$ do not change.
\end{remark}

\begin{corollary}[all computable probability measures]\label{cor:allmeasures}
Let $\mathcal C$ be a countable set of probability measures that contains
$\mu$ and $\lambda$, and with $\nu^+,\nu^-\in\mathcal C$ for every
$\nu\in\mathcal C$. Then the conclusion of Theorem~\ref{thm:measures} holds
for $\mathcal C$ in place of $\cCQ$, with sums over $\nu\in\mathcal C$, and
also with the upper bound $n^{1/4}/8$. This applies to the countable set
$\cCcomp$ of all computable probability measures, and to the set of computable probability
measures with computable rational conditionals at histories of positive mass.
(Paper step; not formally verified.)
\end{corollary}

\begin{proof}
The proof of Theorem~\ref{thm:measures} uses only these properties of the
class, with the fiber sums $\tilde a^\pm_\rho=\sum_{\nu:\,\nu^\pm=\rho}a_\nu$. The
copies $\nu^\pm$ of a probability measure are probability measures: at $1^k$,
$\nu^\pm(1^k0)+\nu^\pm(1^{k+1})=\nu(z_k)+\nu(1^{k+1})=\nu(1^k)$, and inside a
branch the equality comes from $\mu^\pm_{N_j}$. If $\nu$ is computable, then
$\nu^\pm(x)=c_\pm(x)\,\nu(\sigma(x))$ is computable, since $c_\pm$ is rational and
computable. The same formula keeps rational conditionals at histories of
positive mass. Both $\mu$ and $\lambda$ have rational conditionals.
\end{proof}

The proof also works with any $\eta\in(0,2^{-20}]$ in place of $2^{-20}$,
because Theorem~\ref{thm:robust} accepts such $\eta$; then $w^\pm\ge\eta a$
(paper step).

\begin{remark}[semantic complexity weights; proof sketch; not formally verified]\label{rem:exactK}
The book's universal prior gives a model $\nu$ the weight $2^{-K(\nu)}$. On a
set of distinct measures, a natural reading of $K(\nu)$ is
$K_V(\nu):=\min\{K_V(q):q\text{ computes }\nu\}$ for an optimal prefix machine
$V$, and $K_{U_\pm}(\nu)$ is defined in the same way. This differs from the
index convention $K(\nu_i):=K(i)$ of Theorems~\ref{thm:complexity}
and~\ref{thm:exact}. In the index convention, Theorem~\ref{thm:exact} gives
exact weights with a formal proof. The semantic statement below has only a
proof sketch; we do not state it as a theorem. Let
$\mu_*$ output one bit $1$ and then run the source $\mu$. For every optimal
prefix machine $V$ there are two constructed optimal prefix machines $U_-$,
$U_+$ such that the mixtures
\[
\Xi_\pm=\sum_{\nu\in\cCcomp}2^{-K_{U_\pm}(\nu)}\,\nu
\]
over the distinct computable probability measures satisfy
$n^{1/4}/512\le L_n^{\Xi_-}-L_n^{\Xi_+}<n^{1/4}/8$ under $\mu_*$ for
$n\ge3$. The weights are exactly $2^{-K_{U_\pm}(\nu)}$, without normalization;
their ratios are in $\{\frac14,1,4\}$. Each machine reads one of three
prefix-free headers and then simulates $V$. Two headers apply an injective
transform of the model with a uniform inverse, which puts the experiment of
Theorem~\ref{thm:robust} with core weights $\frac15,\frac45$ behind a selector;
the third keeps the model unchanged behind a long header. The two machines
assign the two short headers in opposite order. The exact table of shortest
lengths needs a step that a transformed description is valid only if the
original description is valid. This step (the representation of programs,
validity, and the output convention) is only sketched. The core weights $\frac15,\frac45$ also need a variant of
Theorem~\ref{thm:robust}. Its proof gives this variant with
$\pi_\pm\ge\frac15\mu^+$ in Step~2, and more than $N/2$ good rounds remain.
\end{remark}

\subsection{Remarks}

\paragraph{Consistency with the known part.}
The book's bound $L_n^{\xi}-L_n^{\xi'}\le O\bigl((L_n^{\mu})^{1/2}\bigr)$ allows growth of
order $\sqrt n$ here, since $L_n^\mu\approx n/2$. Our differences grow as
$\log n$ (Theorem~\ref{thm:warmup}) and $n^{1/4}$
(Theorem~\ref{thm:mixture}). Every $\xi_w$ with positive weights is Pareto
optimal~\cite[Sec.~5.3]{hutter2003jmlr}; Pareto optimality compares a
predictor across all environments of the class, so it allows an unbounded
difference between two mixtures in one environment.

\paragraph{Where the answer turns positive.}
The negative answer uses a loss with a discrete decision, such as the 0-1
loss, and a source whose conditionals approach the decision threshold. Three standard cases have a finite limit. Let $\xi,\xi'$ be
Bayes mixtures of probability measures with $\mu$ in the class and weights
$w_\mu,w'_\mu>0$, and let $p_t=\mu(1|x_{<t})$ and $q_t=\xi(1|x_{<t})$. The
facts below are known~\cite[Ch.~3]{hutter2005uai},
\cite{solomonoff1978,hutter2003tit}, \cite[p.~988]{hutter2003jmlr}; we use
them only to mark the boundary of the negative answer.
\begin{itemize}
\item \emph{Log loss.} The expected cumulative log-loss difference is
$D_n(\mu\Vert\xi)-D_n(\mu\Vert\xi')$, where
\[
D_n(\mu\Vert\xi)=\E_\mu\ln\frac{\mu(X_{1:n})}{\xi(X_{1:n})}
\]
is nondecreasing in
$n$ and at most $\ln w_\mu^{-1}$. So the difference converges. If
$w'_\nu/w_\nu\in[c,C]$ for all $\nu$, it stays in $[\ln c,\ln C]$.
\item \emph{Quadratic loss.} The expected excess loss over the informed
predictor at time $t$ is $\E_\mu(p_t-q_t)^2$, and the sum of these terms is at
most $\frac12\ln w_\mu^{-1}$. So each regret converges, and so does the
difference.
\item \emph{A uniform margin.} Assume that there is $\delta>0$ such that
$p_t=\frac12$ or $|2p_t-1|\ge\delta$ at every history of positive mass. A wrong
0-1 decision costs $2d$, with $d=|p_t-\frac12|$, and needs $|p_t-q_t|\ge d$.
So the 0-1 regret of $\Lambda_\xi$ is at most
$\frac4\delta\sum_t\E_\mu(p_t-q_t)^2\le\frac2\delta\ln w_\mu^{-1}$, and the
difference of two regrets converges.
\end{itemize}
Our source has margin $2\gamma_j\to0$ in branch $j$. Its experiment there has
$N_j\gamma_j^2=1/1024$, so the data do not reveal the sign of the bias, and
$N_j\gamma_j=4^j/32$, so the expected gain grows with $j$. The 0-1 loss is discontinuous at the decision boundary, which is
where stability results for continuous losses~\cite{kadanechuang1978} do not
apply. For related recent results on bounded prior ratios with squared and log
loss see~\cite[Thm.~5]{liteng2026}, for a single decision flip
see~\cite[Ex.~5]{suwandi2026}, and see also~\cite[Prop.~4]{dawidvovk2009}.

%======================================================================
\section{Problem 6.3: the conditional-complexity error bound fails}\label{sec:63}
%======================================================================

\subsection{Setting}

In the book's Section~6.2, sequence prediction is posed as an agent problem:
in cycle $k$ the agent guesses the next symbol, and the reward says whether the
guess equals the true symbol $\dot z_k$. With a fixed rule for
ties, the agent \AIxi{} is deterministic, and its whole action-reward history
before cycle $k$ is a function of the true prefix $\dot z_{<k}$. So \AIxi{}
induces a map $\Theta:\Bs\to\B$, where $\Theta(\dot z_{<k})$ is its guess in cycle $k$.
We work with arbitrary maps of this kind.

\begin{definition}
A \emph{deterministic predictor} is any function $\Theta:\Bs\to\B$ (no
computability is assumed). For $x,y\in\Bs$ its number of errors while it
predicts $y$ after $x$ is
\[
E_\Theta(x,y)=\#\{1\le i\le|y|:\ \Theta(x\,y_{<i})\neq y_i\}.
\]
\end{definition}

With $x=\dot z_{<k}$ and $y=\dot z_{k:n}$, $E_\Theta(x,y)$ is the error count
$E_{k:n}$ of the book. The last question of Problem 6.3 asks, under this
error-count reading, whether
\begin{equation}\label{eq:63E}
E_{k:n}\le C\,2^{K(\dot z_{k:n}|\dot z_{<k})}
\end{equation}
holds with a constant $C$ independent of $n$, $k$ and $\dot z$. (The book
writes $\approx$, which we read as ``within a multiplicative constant''. If $C$
may depend on $\dot z$, the book's bound~(6.17) already gives the answer.)

\subsection{The counterexample}

\begin{theorem}\label{thm:63E}
Let $K(y|x)$ be ordinary conditional prefix complexity for an optimal prefix
machine. There is a constant $b$ such that for every deterministic predictor
$\Theta$ and every $d\in\N$ there are $x,y\in\Bs$ with
\[
K(y|x)\le b\qquad\text{and}\qquad E_\Theta(x,y)\ge d-1 .
\]
Consequently, for every deterministic predictor $\Theta$ there is no $C$ with
$E_\Theta(x,y)\le C\,2^{K(y|x)}$ for all $x,y$. In particular, for \AIxi{} with any
fixed tie rule, no $C$ independent of $n,k,\dot z$ satisfies~\eqref{eq:63E}.
(Formal version: every deterministic predictor; the reduction from \AIxi{} to
a deterministic predictor is a paper step.)
\end{theorem}

\begin{proof}
For $m\ge1$ let $\mathrm{bin}(m)$ be the binary representation of $m$ without
leading zeros. Define a computable $F:\Bs\to\Bs$ by $F(\eps)=\eps$ and, for
$|x|=m\ge1$, with $\delta=|\mathrm{bin}(m)|$,
\[
F(x)=0^{2^\delta-m}\,\mathrm{bin}(m).
\]
Let $b=b_F$, so $K(F(x)|x)\le b$ for all $x$. The constant does not depend on
$\Theta$ or $d$.

Fix $\Theta$ and $d$. For $d\le1$ take $x=y=\eps$. For $d\ge2$ define a word
$c=c_1\dots c_d$ by $c_1=1$ and $c_i=1-\Theta(0^{2^d}c_{<i})$ for $2\le i\le d$.
Let $m$ be the number with $\mathrm{bin}(m)=c$; then $2^{d-1}\le m<2^d$ and
$|\mathrm{bin}(m)|=d$. Put
\[
x=0^m,\qquad y=0^{2^d-m}c=F(x),\qquad xy=0^{2^d}c .
\]
Then $K(y|x)\le b$. The bits $c_2,\dots,c_d$ lie in $y$ (since $m<2^d$), and
$\Theta$ guesses each of them wrongly by construction. So $E_\Theta(x,y)\ge d-1$.

For the consequence, given $C$ choose $d$ with $d-1>C\,2^b$; then
\[
E_\Theta(x,y)\ge d-1>C\,2^{b}\ge C\,2^{K(y|x)}.\qedhere
\]
\end{proof}

The proof does not use any property of $\Theta$, so the failure of~\eqref{eq:63E}
is a property of $K(\cdot|x)$: the condition $x=0^m$ hands the program the
number $m$, and the adversarial bits are stored in $m$. This is the length
observation of Chernov and Hutter~\cite{chernovhutter2005}, in the journal
version~\cite[Lemma~5]{chernovhutterschmidhuber2007}: for $x=0^\ell$ they have
$K(\mu|x)=O(1)$ while the one-step log-loss deviation of $M$ is at least
$K(\ell)\ln2$, and they remark after that lemma that $K(y|x)$ can use the length of
the condition. Theorem~\ref{thm:63E} extends this mechanism from one deviation
of $M$ to the error count of every deterministic predictor; we do not claim the
mechanism as new.

\begin{corollary}\label{cor:63rate}
For every deterministic predictor, no bound
$E_{k:n}\le C\,g(k)\,2^{K(\dot z_{k:n}|\dot z_{<k})}$ with $g(k)=o(\log k)$
holds for all $n,k,\dot z$. (Paper step from Theorem~\ref{thm:63E}.)
\end{corollary}

\begin{proof}
In the proof of Theorem~\ref{thm:63E}, $k=m+1\le2^d$, so $E_{k:n}\ge d-1\ge\log_2k-1$
while~$K(y|x)\le\nobreak b$.
\end{proof}

A correction by the length of the condition is natural:
Chernov, Hutter and Schmidhuber note that logarithmic corrections in the
exponent become polynomial factors~\cite[p.~9]{chernovhutterschmidhuber2007}.

\subsection{The second link of the chain, and what stays open}

The ratio form of the question also fails at its second link.

\begin{proposition}\label{prop:63R}
Let $M$ be a universal semimeasure. There is $C_0$ with $K(1|0^\ell)\le C_0$
for all $\ell$, and for every $r\in\N$ some $\ell$ has $2^r M(0^\ell1)<M(0^\ell)$.
So no constant $C$ gives $M(x)/M(xy)\le C\,2^{K(y|x)}$ for all $x,y$.
(Formal version: the universal semimeasure of the library~\cite{milovanovlean}.)
\end{proposition}

\begin{proof}
$K(1|0^\ell)$ is bounded because the output is constant. The strings $0^\ell1$
form a prefix-free set, so $\sum_\ell M(0^\ell1)\le1$ and $M(0^\ell1)\to0$.
$M$ dominates the measure concentrated on $0^\infty$, so $M(0^\ell)\ge c>0$
for all $\ell$. So for every $r$ some $\ell$ has $2^rM(0^\ell1)<c\le M(0^\ell)$.
Given $C$, take $r$ with $2^r\ge C\,2^{C_0}$; then $x=0^\ell$, $y=1$ give
$M(x)/M(xy)>2^r\ge C\,2^{K(y|x)}$.
\end{proof}

This is essentially the example of~\cite[Lemma~5]{chernovhutterschmidhuber2007}.
It is also related to Problem~2.6(ii) of the book~\cite[p.~61]{hutter2005uai}.
Let $\KM(x)=-\log_2M(x)$. In the book's notation, with $y$ the past and $x$ the
continuation, Problem~2.6(ii) asks for the upper bound
\[
\KM(yx)-\KM(y)\le\tilde K(x|y)+K(\ell(y))+O(1),
\]
where $\ell(y)=|y|$, for a
suitable variant $\tilde K$ of conditional complexity, and says that it is open whether an additive
constant independent of $x$ and $y$ is sufficient. This question is settled in
the literature. Hutter's prize page lists Problem~2.6 as solved by Chernov and
Hutter~\cite{hutterprizepage,chernovhutter2005}. For the past $0^\ell$ and the
continuation $1$, $\tilde K(1|0^\ell)=O(1)$ for every variant listed in
Problem~2.6, the monotone $K_*$ included, while
$\KM(0^\ell1)-\KM(0^\ell)\ge K(\ell)-O(1)$ by the proof
of~\cite[Lemma~5]{chernovhutterschmidhuber2007}. So the additive-constant
version of Problem~2.6(ii) fails, and the $K(\ell(y))$ term
of~\cite[Thm.~1]{chernovhutterschmidhuber2007} is tight, as their Lemma~3 states
for Problem~2.6(iii). Proposition~\ref{prop:63R} is a weaker, machine-checked
form of this fact: it shows that the increase is unbounded, while their bound
gives at least $K(\ell)$ for each $\ell$. We make no claim about Problem~2.6.

Theorem~\ref{thm:63E} shows directly that no route gives
$E_{k:n}\lesssim2^{K(\dot z_{k:n}|\dot z_{<k})}$, and
Proposition~\ref{prop:63R} shows that the second link fails on its own. So the
end of Hutter's chain is refuted.
Two questions stay open:
\begin{itemize}
\item the first link for the active agent,
$E_{k:n}^{\AI\xi}\le C\,\xi^{\SP}(\dot z_{<k})/\xi^{\SP}(\dot z_{1:n})$;
\item the bound with Chernov and Hutter's complexity $K_*(y|x*)$, which is
monotone in the condition~\cite{chernovhutterschmidhuber2007}. Our witnesses do
not refute it, since $K_*(y|0^m*)$ is about $K(m)$.
\end{itemize}
Problem 3.13 asks which results remain valid when $K(\nu)$ is replaced by
$K(\nu|x_{<k})$, and states that for $M=M_U$ the posterior weights do not
satisfy $w_\nu(x_{<k})\geqx2^{-K(\nu|x_{<k})}$ \cite[p.~122]{hutter2005uai};
see also~\cite[Sec.~5(h)]{hutter2009open}. Theorem~\ref{thm:63E} is a failure of
posterization of the same kind, for error counts. For related work on
conditional complexity see also~\cite{andreevposobinshen2017}.

%======================================================================
\section{Problem 2.7: the normalized universal predictor}\label{sec:27}
%======================================================================

\subsection{Setting and the known part}\label{sec:27setting}

For a universal semimeasure $M$, the normalized predictor is
\[
q_M(b|x)=\frac{M(xb)}{M(x0)+M(x1)},
\]
which is the book's $\Mnorm(x_t|x_{<t})$ with $x=x_{<t}$. Let $Q_M$ be the
probability measure with conditionals $q_M$. For a computable probability
measure $\mu$ and $t\ge1$ define the expected one-step $L_1$ and squared errors
\[
A_t=\sum_{|x|=t-1}\sum_{b\in\B}\bigl|\mu(x)\,q_M(b|x)-\mu(xb)\bigr|,\qquad
S_t=\sum_{|x|=t-1}\mu(x)\sum_{b\in\B}\bigl(q_M(b|x)-\mu(b|x)\bigr)^2,
\]
where $\mu$-null histories contribute $0$. For $\mu(x)>0$ let
\[
a_t(x)=\sum_{b\in\B}\bigl|q_M(b|x)-\mu(b|x)\bigr|,\qquad
s_t(x)=\sum_{b\in\B}\bigl(q_M(b|x)-\mu(b|x)\bigr)^2 .
\]
So $A_t=\E_\mu a_t(X_{<t})$ and $S_t=\E_\mu s_t(X_{<t})$. For binary distributions the
squared error is half the square of the $L_1$ error, and Jensen's inequality
gives
\begin{equation}\label{eq:AS}
S_t\ge\tfrac12A_t^2 .
\end{equation}
Recall $\kappa_t=2^{-K(t)}$.

For deterministic environments the answer to the first question is known.

\begin{proposition}[known; Hutter 2007]\label{prop:27F}
For every universal $M$ there is $c_M>0$ with $q_M(b|x)\ge c_M\,2^{-K(|x|+1)}$
for all $x\in\Bs$ and $b\in\B$. In particular, along any sequence the
normalized predictor gives probability at least $c_M2^{-K(t)}$ to the wrong
bit at every time $t$. (Formal version: the universal semimeasure of the
library~\cite{milovanovlean}.)
\end{proposition}

\begin{proof}
For a Solomonoff prior $V$, Hutter~\cite[eq.~(17)]{hutter2007bayesconf} proves
$2^{-K(n)}\leqx V(\bar x_n|x_{<n})$ for all $n$, all $x_{1:n}$ and
$\bar x_n\neq x_n$; since $x_n$ is arbitrary, this covers every history and
every next bit. $V$ and $M$ are both universal and both lsc, so each
dominates the other: $aV\le M\le A'V$ for constants $a,A'>0$. Then
$q_M(b|x)\ge M(xb)/M(x)\ge(a/A')\,V(xb)/V(x)$.
\end{proof}

For a computable sequence $z$ the wrong-bit probability is also
$O(2^{-K(t)})$: the strings $z_{<t}\bar z_t$ form a c.e.\ prefix-free set, so
$g(t)=M(z_{<t}\bar z_t)$ is lsc with $\sum_tg(t)\le1$, hence $g(t)\le C_g\kappa_t$
by~\eqref{eq:coding}, while $M(z_{1:t})\ge u_z>0$; so
$q_M(\bar z_t|z_{<t})\le g(t)/M(z_{1:t})\le(C_g/u_z)\kappa_t$. This two-sided rate of
order $2^{-K(t)}$ is consistent with Leike and Hutter~\cite{leikehutter2015nicod}.
The prize page states that a solution of part ``2.7u'' was written up in 2005;
it is not public~\cite{hutterprizepage}. Hutter's open-problems
paper~\cite{hutter2009open} restates the question for $\Mnorm$ and conjectures
a positive answer. The rest of this section treats stochastic $\mu$.

\subsection{Tools}

\begin{lemma}[normalization keeps dominance]\label{lem:normdom}
$Q_M(x)\ge M(x)/M(\eps)\ge M(x)$ for all $x$.
\end{lemma}

\begin{proof}
\[
Q_M(x)=\prod_{i=1}^{|x|}\frac{M(x_{\le i})}{M(x_{<i}0)+M(x_{<i}1)}
\ge\prod_{i=1}^{|x|}\frac{M(x_{\le i})}{M(x_{<i})}=\frac{M(x)}{M(\eps)}.
\qedhere
\]
\end{proof}

The next lemma forces a discrepancy at prescribed times with arbitrary gaps
between them.

\begin{lemma}[selective forcing]\label{lem:forcing}
Let $P$ and $Q$ be probability laws on $\B^N$ with conditionals $p,q$, with
$q(b|x)>0$ everywhere. Let $F\subseteq\{1,\dots,N\}$ with $|F|=m$, and for
each $t\in F$ and $|x|=t-1$ let $f_t(x)\in\B$ satisfy $p(f_t(x)|x)\le\rho<1$.
Let $\nu$ be the law that follows $p$ at times outside $F$ and outputs $f_t(x)$
at times $t\in F$. If $0<r\le1$ and $Q(z)\ge r\,\nu(z)$ for all $|z|=N$, then
\[
\sum_{t\in F}A_t(P,Q)\ge2(r-\rho^m),\qquad
A_t(P,Q):=\E_P\sum_b|q(b|X_{<t})-p(b|X_{<t})| .
\]
\end{lemma}

\begin{proof}
Let $\widehat Q$ be the law with conditionals $q$ at times in $F$ and $p$ at
other times, and let $E$ be the event that $X_t=f_t(X_{<t})$ for all $t\in F$.
By successive conditioning, $P(E)\le\rho^m$.

We show $\widehat Q(E)\ge r$. For a prefix $x$ let $h(x)$ be the
$\widehat Q$-probability of $E$ given $x$. We prove by backward induction on
$|x|$ that $h(x)\ge r\,\nu(x)/Q(x)$. This is trivial if $\nu(x)=0$. For
$|x|=N$ and $\nu(x)>0$, $h(x)=1$ and the claim is the domination
hypothesis. If $|x|+1\in F$ and $f=f_{|x|+1}(x)$, then $\nu(xf)=\nu(x)$,
$Q(xf)=Q(x)q(f|x)$ and
$h(x)=q(f|x)h(xf)\ge q(f|x)\,r\nu(x)/(Q(x)q(f|x))=\nobreak r\nu(x)/Q(x)$.
If $|x|+1\notin F$, then $\nu(xb)=\nu(x)p(b|x)$ and $Q(xb)=Q(x)q(b|x)$, so
\[
h(x)=\sum_bp(b|x)h(xb)\ge r\frac{\nu(x)}{Q(x)}\sum_b\frac{p(b|x)^2}{q(b|x)}
\ge r\frac{\nu(x)}{Q(x)}
\]
by~\eqref{eq:chi2}. At the root, $\widehat Q(E)=h(\eps)\ge r$.

Hence $\|P-\widehat Q\|_1\ge2(\widehat Q(E)-P(E))\ge2(r-\rho^m)$. By the hybrid
inequality~\eqref{eq:hybrid} with the conditionals of $\widehat Q$, which equal
those of $P$ outside $F$, $\|P-\widehat Q\|_1\le\sum_{t\in F}A_t(P,Q)$.
\end{proof}

\begin{lemma}[changing the environment]\label{lem:envchange}
Fix $q$. For probability laws $P,R$ and
\begin{gather*}
A_t(P,Q)=\sum_{|x|=t-1}\sum_b|P(x)q(b|x)-P(xb)|,\\
|A_t(P,Q)-A_t(R,Q)|\le\|P_{t-1}-R_{t-1}\|_1+\|P_t-R_t\|_1,
\end{gather*}
where $P_t$ is the marginal on $\B^t$. If $R=(1-\delta)P+\delta L$ for a
probability law $L$, the right side is at most $4\delta$.
\end{lemma}

\begin{proof}
This $A_t(P,Q)$ is the one of Lemma~\ref{lem:forcing}: if $P(x)>0$ then, with
$p(b|x)=P(xb)/P(x)$, $|P(x)q(b|x)-P(xb)|=P(x)\,|q(b|x)-p(b|x)|$, and $P$-null histories contribute
$0$ to both. The present form needs no conditionals of $P$.
Apply the triangle inequality termwise and use $\sum_bq(b|x)=1$. For the
mixture, $\|P_t-R_t\|_1=\delta\|P_t-L_t\|_1\le2\delta$.
\end{proof}

\subsection{Lower bounds along computable selections}

\begin{theorem}[selections]\label{thm:selection}
Let $M$ be a universal semimeasure, $\mu$ a computable probability measure,
and $d_1<d_2<\cdots$ a computable increasing sequence of times. Let
$\ell_n=6\lceil\log_2(n+2)\rceil$, $B_0=0$, $B_n=\sum_{k=1}^n\ell_k$, and
$F_n=\{d_j:B_{n-1}<j\le B_n\}$. There are $c>0$ and $n_0$ such that for every
$n\ge n_0$ some $t\in F_n$ has
\[
A_t\ge c\,\kappa_t\qquad\text{and}\qquad S_t\ge\frac{c^2}2\,\kappa_t^2 .
\]
The constants depend on $M$, $\mu$, the machine defining $K$, and the
sequence $d$. The formal version assumes that $\mu$ has rational conditionals
given by a computable function.
\end{theorem}

\begin{proof}
\emph{Alternatives.} Let $\tilde\mu_n=(1-2^{-n})\mu+2^{-n}\lambda$. These measures
have full support, and their conditionals are computable uniformly in $n$,
because a cylinder of length $k$ has $\tilde\mu_n$-mass at least $2^{-n-k}$. At each
$t\in F_n$ and history $x$ of length $t-1$, one bit has $\tilde\mu_n$-conditional
probability at most $\frac12$; a search over rational approximations finds a
bit $f_t(x)$ with $\tilde\mu_n(f_t(x)|x)<\frac23$. Let $\tilde\nu_n$ follow $\tilde\mu_n$ outside
$F_n$ and output $f_t(x)$ at $t\in F_n$. The family $(\tilde\nu_n)$ is uniformly
computable.

\emph{One domination constant.} $R=\sum_{n\ge1}\kappa_n\tilde\nu_n$ is an lsc
semimeasure. By universality and Lemma~\ref{lem:normdom} there is $a\in(0,1]$
with $Q_M\ge M\ge aR\ge a\,\kappa_n\tilde\nu_n$ for every $n$.

\emph{Forcing.} Apply Lemma~\ref{lem:forcing} with $P=\tilde\mu_n$, $Q=Q_M$, $\nu=\tilde\nu_n$,
$N=\max F_n$, $\rho=\frac23$, $m=\ell_n$ and $r=a\kappa_n$:
$\sum_{t\in F_n}A_t(\tilde\mu_n,Q_M)\ge2a\kappa_n-2(2/3)^{\ell_n}$. By
Lemma~\ref{lem:envchange},
\[
\sum_{t\in F_n}A_t\ge2a\kappa_n-2(2/3)^{\ell_n}-4\ell_n2^{-n}.
\]
Now $(2/3)^{\ell_n}\le(n+2)^{-6\log_2(3/2)}$ with $6\log_2(3/2)>3$, and
$\kappa_n\ge c_0(n+1)^{-2}$ by~\eqref{eq:wlower}. So for $n\ge n_0$,
$\sum_{t\in F_n}A_t\ge a\,\kappa_n$.

\emph{Coding.} The sets $F_n$ are finite, pairwise disjoint and computable
uniformly in $n$. So $g(n)=\sum_{t\in F_n}\kappa_t$ is lsc and
$\sum_ng(n)\le\sum_t\kappa_t\le1$. By~\eqref{eq:coding},
$\sum_{t\in F_n}\kappa_t\le C_g\,\kappa_n$. Hence
$\sum_{t\in F_n}A_t\ge(a/C_g)\sum_{t\in F_n}\kappa_t$, and some $t\in F_n$ has
$A_t\ge(a/C_g)\kappa_t$. Put $c=a/C_g$ and use~\eqref{eq:AS}.
\end{proof}

The constant $c$ depends on the selection. A union over all selections
therefore does not give a bound at every large $t$.

\begin{corollary}[logarithmic windows]\label{cor:window}
For every universal $M$ and computable $\mu$ there are $c>0$, $C_1$ and $T$ such
that every window $\{u,\dots,u+C_1(\lfloor\log_2(u+2)\rfloor+1)\}$ with $u\ge T$
contains a time $t$ with $A_t\ge c\,\kappa_t$ and $S_t\ge\frac{c^2}2\kappa_t^2$. Hence
there are $\Omega(N/\log N)$ such times up to $N$. (Formal version: the
window statement, with rational conditionals; the count is a paper step.)
\end{corollary}

\begin{proof}
Take $d_j=j$. If $B_{n-1}<u\le B_n$, the block $F_{n+1}$ lies in
$[u,u+\ell_n+\ell_{n+1}]$, and $n\le u$, so
$\ell_n+\ell_{n+1}=O(\log(u+2))$. The count follows from
$B_n=O(n\log n)$.
\end{proof}

\begin{corollary}[dyadic blocks]\label{cor:dyadic}
For every universal $M$ and computable $\mu$ there are $c>0$ and $n_0$ such that
for every $n\ge n_0$ some $t\in(2^n,2^{n+1}]$ has $A_t\ge c\,\kappa_t$ and
$S_t\ge\frac{c^2}2\kappa_t^2$. (Formal version: rational conditionals.)
\end{corollary}

\begin{proof}
Apply Corollary~\ref{cor:window} with $u=2^n+1$. For large $n$,
\[
2^n+1+C_1(\lfloor\log_2(2^n+3)\rfloor+1)\le2^{n+1}.\qedhere
\]
\end{proof}

The formal proof of Corollary~\ref{cor:dyadic} follows an earlier, direct
route: in each dyadic block it compares $\mu$ with alternatives that follow
$\mu$ for $2^n$ steps and then a computable path of small $\mu$-probability,
and uses the hybrid inequality~\eqref{eq:hybrid} and Markov's inequality.

\begin{corollary}[c.e.\ sets of times]\label{cor:ce}
For every universal $M$, computable $\mu$ and infinite computably enumerable set
$D$ of times, there is $c>0$ such that infinitely many $t\in D$ have
$A_t\ge c\,\kappa_t$ and $S_t\ge\frac{c^2}2\kappa_t^2$. (Paper proof from
Theorem~\ref{thm:selection}; not formalized.)
\end{corollary}

\begin{proof}
An infinite c.e.\ set contains a computable increasing sequence: after $d_j$,
wait for an enumerated element larger than $d_j$. Apply
Theorem~\ref{thm:selection}.
\end{proof}

\subsection{No general upper bound by a power of $2^{-K(t)}$}

The second question of Problem 2.7 asks for which sequences the error is
$\eqx2^{-K(t)}$. For stochastic environments, an upper bound of this order does
not hold for all universal $M$, even with any positive power.

\begin{theorem}[spikes for the fair coin]\label{thm:spikes}
Fix an optimal prefix machine defining $K$ and $\kappa_t=2^{-K(t)}$. For every
universal semimeasure $V$ with $V(\eps)=1$ there is an lsc semimeasure $\tau$
with $\tau(\eps)=1$ such that $M=\frac12V+\frac12\tau$ is universal and, for the
fair coin $\mu=\lambda$,
\[
\forall\alpha>0\ \forall C\ \forall T\ \exists t\ge T:\qquad
S_t>C\,\kappa_t^{\alpha}\quad\text{and}\quad A_t>C\,\kappa_t^{\alpha}.
\]
In particular, there is a universal $M$ whose fair-coin expected errors are not
$O(2^{-\alpha K(t)})$ for any $\alpha>0$.
\end{theorem}

\begin{proof}
\emph{Busy-beaver times.} Let
$\beta(n)=\max(\{0\}\cup\{j:K(j)\le n\})$, with $K$ from the same machine as
$\kappa$. It is finite, nondecreasing, unbounded and lsc. Let $t_n=\beta(n)+1$;
then $K(t_n)>n$, so $\kappa_{t_n}<2^{-n}$. Let
$\gamma_n=2^{-(2\lfloor\log_2n\rfloor+1)}$ for $n\ge1$; then
$\sum_n\gamma_n=1$ and $\gamma_n\ge1/(2n^2)$.

\emph{Stopping semimeasures.} For $L\ge0$ and $i\in\B$ let
$\tau_{L,i}(x)=\lambda(x)$ if $|x|\le L$, or if $|x|=L+1$ and the last bit of $x$
is $i$; and $\tau_{L,i}(x)=0$ otherwise. This is a semimeasure, and
$\tau_{L,i}(x)$ is nondecreasing in $L$. So
$\tau_i=\sum_n\gamma_n\tau_{\beta(n),i}$ is lsc, with $\tau_i(\eps)=1$. Let
$M^{(i)}=\frac12V+\frac12\tau_i$; it is lsc, has root $1$, and is universal.

\emph{The two predictors differ at the stopping times.} Fix $t$ and $|x|=t-1$.
Let $W_t=\sum_{n:\beta(n)=t-1}\gamma_n$ and
$h_t=\sum_{n:\beta(n)\ge t}\gamma_n+\frac12W_t$. Summing over $n$ by cases,
$\tau_i(x0)+\tau_i(x1)=h_t\lambda(x)$ for both $i$, so the two normalizers agree:
\[
D_t(x):=M^{(i)}(x0)+M^{(i)}(x1)=\tfrac12\bigl(V(x0)+V(x1)+h_t\lambda(x)\bigr).
\]
Also $M^{(0)}(x0)-M^{(1)}(x0)=\frac12W_t\lambda(x0)$, so
\[
q_{M^{(0)}}(0|x)-q_{M^{(1)}}(0|x)=\frac{W_t\lambda(x)}{4D_t(x)} .
\]
Since $\sum_{|x|=t-1}D_t(x)\le1$, Cauchy--Schwarz gives
$\E_\lambda[\lambda(X)/D_t(X)]=\sum_x\lambda(x)^2/D_t(x)\ge1$, and Jensen gives
$\E_\lambda[(q_{M^{(0)}}-q_{M^{(1)}})^2]\ge W_t^2/16$. For the fair coin,
$s_t(x)=2(q_M(0|x)-\frac12)^2$ with $s_t$ as in Section~\ref{sec:27setting}, and $2(u-\frac12)^2+2(v-\frac12)^2\ge(u-v)^2$.
Hence
\[
S_t(M^{(0)})+S_t(M^{(1)})\ge\frac{W_t^2}{16}.
\]
\emph{Pigeonhole.} At $t=t_n$, $W_{t_n}\ge\gamma_n$. So for each $n$ some
$i$ has $S_{t_n}(M^{(i)})\ge\gamma_n^2/32$, and one fixed $i$ does so for
infinitely many $n$. Take $\tau=\tau_i$. Along these $n$, $t_n\to\infty$ and
\[
\frac{S_{t_n}}{\kappa_{t_n}^\alpha}\ge\frac{\gamma_n^2}{32}2^{\alpha n}\ge\frac{2^{\alpha n}}{128\,n^4}\to\infty .
\]
For the fair coin, $a_t(x)=2|q_M(0|x)-\frac12|\ge s_t(x)$, so $A_t\ge S_t$.
\end{proof}

\begin{remark}[Solomonoff priors]
If $V$ is the Solomonoff prior of a universal monotone machine, then so is $M$:
a machine reads one bit and then either simulates the machine of $V$ or samples
from $\tau$. With universality by adjunction this is the definition of Wood,
Sunehag and Hutter~\cite[Def.~2]{woodsunehaghutter2013}; see also their Lemma~13
and Theorem~14. This step is a paper argument; Lean checks the statement for
universal lsc semimeasures. Theorem~\ref{thm:spikes} is about one $M$; it does
not say that every universal $M$ has such spikes.
\end{remark}

\subsection{What remains open}

The first question in its strongest stochastic form is:
\begin{equation}\label{eq:E}\tag{E}
\forall M\ \forall\mu\ \exists c>0\ \exists T\ \forall t\ge T:\quad
A_t\ge c\,\kappa_t\quad(\text{or } S_t\ge c\,\kappa_t^2),
\end{equation}
for universal $M$ and computable $\mu$. We do not know whether (E) holds,
also for the fair coin and for Solomonoff priors. The results above restrict a
counterexample.

\begin{proposition}[quiet sets are immune]\label{prop:immune}
Let $M$ be universal and $\mu$ computable. Call an infinite set $D$ of times
\emph{quiet} if $A_t/\kappa_t\to0$ as $t\to\infty$ in $D$, or
$S_t/\kappa_t^2\to0$ as $t\to\infty$ in $D$. Then a quiet set has no infinite
c.e.\ subset. If \textup{(E)} fails for $M,\mu$, a quiet set exists. Moreover, with
$c$, $C_1$ and $T$ from Corollary~\ref{cor:window}, every window
$\{u,\dots,u+C_1(\lfloor\log_2(u+2)\rfloor+1)\}$ with $u\ge T$ contains a time
$t$ with $A_t\ge c\,\kappa_t$ and $S_t\ge\frac{c^2}2\kappa_t^2$. So from some point on, a quiet set omits at
least one time in each such window.
\end{proposition}

\begin{proof}
An infinite c.e.\ subset of a quiet set $D$ would contain infinitely many $t$
with $A_t\ge c\,\kappa_t$ and $S_t\ge\frac{c^2}2\kappa_t^2$ (Corollary~\ref{cor:ce}),
contrary to $A_t/\kappa_t\to0$ or $S_t/\kappa_t^2\to0$ along $D$. If (E) fails, then
$\liminf_tA_t/\kappa_t=0$ (or $\liminf_tS_t/\kappa_t^2=0$), which gives a quiet set.
The last claim is Corollary~\ref{cor:window}.
\end{proof}

Two simpler obstacles are known to us. Exact symmetry, $A_t=0$ for the fair
coin at infinitely many $t$, is impossible: it would force
$Q_M(0^N)\to0$, while $Q_M(0^N)\ge M(0^N)\ge c>0$. So a counterexample needs
small nonzero errors on a quiet set. And a summed bound does not help
directly: if $M\ge c_\mu\mu$, then $\sum_tS_t\le\ln(1/c_\mu)$ by the standard KL
argument~\cite{hutter2003tit}, which allows spikes as in
Theorem~\ref{thm:spikes}.

%======================================================================
\section{Formal verification}\label{sec:formal}
%======================================================================

\paragraph{Toolchain.}
Lean~4.35.0~\cite{lean4}, Mathlib at tag \texttt{v4.35.0}~\cite{mathlib2020},
and \texttt{kolmogorov-complexity-lean}~\cite{milovanovlean} at commit
\texttt{17c50e560c337a8a702f8850a711f80a56b3cf91}. A check script builds every
file in the directory \path{HutterPrizes/} of the repository, fails on any
\texttt{sorry} in these files, and prints the axioms of every theorem in
Table~\ref{tab:lean}; it fails on any axiom other than \texttt{propext},
\texttt{Classical.choice} and \texttt{Quot.sound}. Every theorem in
Table~\ref{tab:lean} passes this check.
The repository is at \repourl.

\paragraph{Process.}
For the results on universal mixtures (Theorems~\ref{thm:mixture},
\ref{thm:complexity} and~\ref{thm:exact}), on mixtures over measures
(Theorem~\ref{thm:measures}), Problem 6.3 (Theorem~\ref{thm:63E},
Proposition~\ref{prop:63R}) and Problem 2.7 (Proposition~\ref{prop:27F},
Theorem~\ref{thm:selection}, Corollaries~\ref{cor:window}
and~\ref{cor:dyadic}, Theorem~\ref{thm:spikes}), a Lean statement was frozen
before the proof was completed, with a note that lists its definitions, its
quantifiers and the dependence of each constant. A fresh-context AI reviewer
(Claude) compared each statement with the book's text and recorded the claim it
accepts. Two targets were changed after their reviews: in one, a hypothesis
that excluded valid selections was removed; the target of
Theorem~\ref{thm:exact} gained the clauses that its reviewer asked for (the
per-index weight ratio, universality, and the renumbering of the enumeration).
The final statements, including those of Theorems~\ref{thm:complexity},
\ref{thm:exact} and~\ref{thm:measures}, are identical to the frozen ones. Theorem~\ref{thm:warmup} was reviewed in the same way after its
proof. Theorems~\ref{thm:robust} and~\ref{thm:aggregate} are intermediate
results; we compared their statements with the paper proofs.

\paragraph{Conventions in the formal statements.}
Strings are \lean{List Bool}. Semimeasures take values in $[0,\infty]$; the
library class \lean{IsLowerSemicomputableContinuousSemimeasure} is our $\cMone$
(lsc, root mass exactly $1$). \lean{P33Effective.IsUniversal} means
domination of every member of this class; by Lemma~\ref{lem:root} this is
equivalent to the book's universality. $K(y|x)$ is \lean{KP U y x} and $K(t)$ is
\lean{KPNat U t} for a machine \lean{U} with
\lean{IsOptimalPrefixConditional U}. In Problem 2.7, the true $\mu$ is given by a
computable \lean{r : List Bool -> Q} with $\mu(1|x)=r(x)$ (rational
conditionals); the error sums \lean{expAbsErr} and \lean{expSqErr} are $A_t$
and $S_t$. In Problem 3.3, the predictor breaks ties toward~1, except in
\lean{MixtureCounterexample.counterexample}, which breaks them toward~0, and
\lean{transport} identifies histories with strings. For
Theorem~\ref{thm:measures}, \lean{RatMeasure} is the index set of $\cCQ$
(computable \lean{r : List Bool -> Q} with values in $[0,1]$),
\lean{measureMass} is $\nu_r$, and \lean{classMixture} applied to $w$ is
$\sum_rw_r\nu_r$.
The lemmas \lean{ratMeasure_countable} and \lean{exists_positive_prior} in the
same file show that the index set is countable and has a positive prior, so
the hypothesis on $a$ can be met.

\begin{table}[p]
\centering\small\setlength{\defaultaddspace}{0.4em}
\begin{tabularx}{\textwidth}{@{}>{\raggedright\arraybackslash}p{2.7cm}>{\raggedright\arraybackslash}X>{\raggedright\arraybackslash}p{4.0cm}@{}}
\toprule
Paper & Lean name & Scope of the formal version \\
\midrule
Thm.~\ref{thm:warmup} & \lean{MixtureCounterexample.counterexample} & lower bound and unboundedness; upper bound by paper step; ties toward 0 \\
\addlinespace
Thm.~\ref{thm:robust} & \lean{RobustExperiment.theorem3} & as stated \\
\addlinespace
Thm.~\ref{thm:aggregate} & \lean{RobustExperiment.Theorem4.theorem4_statement}, \lean{RobustExperiment.Theorem4.xi_dominates}, \lean{RobustExperiment.Theorem4.xi_semimeasure} & as stated \\
\addlinespace
Thm.~\ref{thm:mixture} & \lean{P33Mixture.mixture_counterexample} & library enumeration of $\cMone$; universality by paper step \\
\addlinespace
Thm.~\ref{thm:aggregate} with a universal background & \lean{P33Effective.effective_counterexample}, \lean{P33Effective.loss_gap_tendsto_atTop} & two universal $\xi_\pm\in\cMone$; a different pair from Thm.~\ref{thm:mixture} \\
\addlinespace
Thm.~\ref{thm:complexity} & \lean{P33Complexity.complexity_counterexample} & library enumeration of $\cMone$; universality by paper step \\
\addlinespace
Thm.~\ref{thm:exact} & \lean{P33Exact.exact_counterexample} & $\tau$ renumbers the library enumeration; bound $n^{1/4}/512$, $n\ge3$ \\
\addlinespace
Thm.~\ref{thm:measures} & \lean{P33Measures.measure_class_counterexample} & class indexed by total computable rational $r$; $\eta=2^{-20}$; upper bound by paper step \\
\addlinespace
Known part of 3.3 (squared distance) & \lean{HutterPrizes.two_semimeasures_sum_sq_le} & calibration \\
\addlinespace
Thm.~\ref{thm:63E} & \lean{HutterPrizes.errorcount_complexity_bound_fails}, \lean{HutterPrizes.no_errorcount_complexity_bound} & every deterministic predictor $\Theta$; the reduction from \AIxi{} to $\Theta$ by paper step \\
\addlinespace
Prop.~\ref{prop:63R} & \lean{HutterPrizes.uniform_conditional_ratio_fails} & library's universal $M$ only \\
\addlinespace
Prop.~\ref{prop:27F} & \lean{HutterPrizes.normalized_lower_bound} & library's universal $M$ only \\
\addlinespace
Lemma~\ref{lem:forcing} & \lean{P27.selective_forcing}, \lean{P27.selective_forcing_bound} & finite form \\
\addlinespace
Thm.~\ref{thm:selection} & \lean{P27Selection.selection_lower_bound} & rational conditionals \\
\addlinespace
Cor.~\ref{cor:window} & \lean{P27Selection.log_window_lower_bound} & rational conditionals \\
\addlinespace
Cor.~\ref{cor:dyadic} & \lean{P27Stochastic.expected_lower_bound_io} (with \lean{P27.block_bound}, \lean{P27.block_bound_of_joint}, \lean{P27.hybrid_inequality}) & rational conditionals \\
\addlinespace
Thm.~\ref{thm:spikes} & \lean{P27Spikes.no_power_upper_bound_mix}, \lean{P27Spikes.no_power_upper_bound} & as stated; Solomonoff-prior remark by paper step \\
\bottomrule
\end{tabularx}
\caption{Paper results and their Lean theorems.}
\label{tab:lean}
\end{table}

\paragraph{What the formal proofs do not cover.}
Lemma~\ref{lem:root} and Remark~\ref{rem:MU} (the passage from $\cMone$ to the
book's class); the upper bound $H_n/16$ in Theorem~\ref{thm:warmup}; the
universality of the mixtures in Theorems~\ref{thm:mixture}
and~\ref{thm:complexity} (the formal statements give lower semicomputability
and the weights; the separate formal result
\lean{P33Effective.effective_counterexample} gives two universal
semimeasures); the computability of the weights of Theorem~\ref{thm:mixture};
the constant $1/256$ and the case $n=2$ in Theorem~\ref{thm:exact} (the
formal bound is $n^{1/4}/512$ for $n\ge3$);
the upper bound in Theorem~\ref{thm:measures}, Remark~\ref{rem:indices} (the
passage from indices to distinct measures), Corollary~\ref{cor:allmeasures}
(all computable probability measures) and other values of $\eta$;
Remark~\ref{rem:exactK}, which has only a proof sketch; the known positive
cases of Section~\ref{sec:33} (log loss, quadratic loss, a uniform margin);
the reduction from \AIxi{} to a deterministic predictor in
Theorem~\ref{thm:63E}; Propositions~\ref{prop:63R} and~\ref{prop:27F} for
every universal $M$ (the formal statements are for the library's $M$);
Theorem~\ref{thm:selection} and its
corollaries for computable $\mu$ with irrational conditionals;
the count in Corollary~\ref{cor:window}; Corollary~\ref{cor:ce} and Proposition~\ref{prop:immune}; Corollary~\ref{cor:63rate};
and the Solomonoff-prior remark after Theorem~\ref{thm:spikes}.

%======================================================================
\section{Related work}\label{sec:related}
%======================================================================

\paragraph{Problem 3.3.}
We found no earlier answer. The closest antecedents are the book's Section~3.6
and Problem~3.6~\cite[pp.~98, 120]{hutter2005uai} and
Hutter~\cite[p.~988]{hutter2003jmlr}: with $\cM=\{\mu_1,\mu_2\}$ and
$\mu(1|x_{<t})=\frac12\pm\varepsilon_t$, the regret
$E_n^{\Theta_\xi}-E_n^{\Theta_\mu}$ of one mixture diverges. They do not compare
two mixtures. Theorem~\ref{thm:warmup} is in the same family; the selector
construction of Theorems~\ref{thm:aggregate} and~\ref{thm:mixture}, which makes
the gap survive a universal background, is new as far as we know, and so is
its extension to mixtures over computable probability measures with
priors $w^\pm\ge2^{-20}a$ built from any positive base prior $a$ (Theorem~\ref{thm:measures}),
and to two universal mixtures with exact weights $2^{-K(i)}$ that differ only
in the reference machine (Theorem~\ref{thm:exact}). The 2024
book~\cite[Exercise~3.10]{hutterquarelcatt2024} keeps only the known part of
Problem 3.3. Hutter's open-problems list~\cite{hutter2009open} does not list the
two-mixture loss question. For universal mixtures and Solomonoff priors we use
the class relations of Wood, Sunehag and Hutter~\cite{woodsunehaghutter2013}.
Leike and Hutter~\cite{leikehutter2015bad} study bad universal priors in the
agent setting; they exclude horizon~1.
Lattimore and Hutter~\cite{lattimorehutter2015} show that convergence of
universal mixtures can fail on individual Martin-L\"of random sequences.

\paragraph{Problem 6.3.}
The mechanism of Theorem~\ref{thm:63E} is due to Chernov and
Hutter~\cite{chernovhutter2005,chernovhutterschmidhuber2007}; our contribution
is the extension to error counts of every deterministic predictor and the
resulting negative answer. Their results with the complexity $K_*$ and with a
length correction remain the positive side of the problem.

\paragraph{Problem 2.7.}
Proposition~\ref{prop:27F} is known~\cite{hutter2007bayesconf}. Wood, Sunehag
and Hutter~\cite[Lemma~15]{woodsunehaghutter2013} give the defect floor
$2^{-K(t)}$ for Solomonoff priors. Leike and Hutter~\cite{leikehutter2015nicod}
give the deterministic rate. Bauwens and Kuznetsov~\cite{bauwenskuznetsov2026}
study $\Mnorm$ for all maximal lsc semimeasures at positions forced by a
decidable tree, with pointwise upper bounds on sums of powers of errors on every
branch and a lower bound on one constructed branch. Their setting has no
stochastic $\mu$ and no expectations; where the settings overlap ($\mu$ on a
decidable tree), their upper bound and our Corollary~\ref{cor:ce} are consistent, and
neither implies the other. We did not find
Theorems~\ref{thm:selection}--\ref{thm:spikes} in the literature. The prediction
of selected bits by $\Mnorm$ (Problem 2.9) was solved by Lattimore, Hutter and
Gavane~\cite{lattimorehuttergavane2011}.

%======================================================================
\section{Conclusion}
%======================================================================

Problem 3.3 and the last question of Problem 6.3 have negative answers.
For Problem 3.3, the gap appears in the book's own setting of a countable
class of computable probability measures, with priors built from any positive
base prior and with bounded prior ratios. It also appears with full universal classes of
semimeasures, and for two universal priors with the exact weights $2^{-K(i)}$
for two constructed reference machines and one constructed enumeration. Solomonoff's mixtures $M_U$ of two
universal monotone machines are not covered here. The gap is a property
of the 0-1 loss near the decision boundary: for log loss, quadratic loss and
sources with a uniform margin the difference is known to stay bounded. For Problem 2.7, the normalized predictor has expected errors of
order $2^{-K(t)}$ infinitely often along every computable selection of times
(with a constant that depends on the selection), and no
upper bound by a power of $2^{-K(t)}$ holds for all universal priors. Three
questions stay open: the every-large-$t$ lower bound (E) for $\Mnorm$; the
numerator improvement $E_{k:n}\lesssim\xi^{\SP}(\dot z_{<k})/\xi^{\SP}(\dot z_{1:n})$
for the active agent; and the bound of Problem 6.3 with the monotone conditional
complexity $K_*$.

\section*{Acknowledgments}
\paragraph{Use of AI tools.}
This work was carried out with substantial help from large language models.
ChatGPT Pro (OpenAI) proposed the proof sketches and the Lean code. Claude
(Anthropic) reviewed each formal statement against the problem text, refereed
the paper proofs, searched the literature, and drafted this manuscript. Every
theorem marked as formally verified was checked by the Lean~4 kernel with
Mathlib and the \texttt{kolmogorov-complexity-lean} library; the statements,
proofs and review notes are available at \repourl. The author is not a
specialist in algorithmic information theory; the author directed the work and
chose the problems and the checks. The correctness claims rest on the Lean
verification and on the independent model reviews described in
Section~\ref{sec:formal}. The author takes full responsibility for the content,
including any errors.

\bibliographystyle{plainurl}
\setlength{\bibsep}{2pt plus 1pt}
{\raggedright
\bibliography{refs}
\par}

\end{document}
